\documentclass[11pt]{amsart}
\usepackage[a4paper,margin=25mm]{geometry}
\usepackage{amssymb,amsmath,mathtools,mathrsfs}
\newcommand{\eps}{\varepsilon}
\newcommand{\dd}{\,\mathrm d}
\usepackage{xcolor}
\usepackage[colorlinks=true,linkcolor=blue!50!black,urlcolor=blue!50!black]{hyperref}
\newtheorem{theorem}{Theorem}[section]
\newtheorem{lemma}[theorem]{Lemma}
\newtheorem{proposition}[theorem]{Proposition}
\newtheorem{corollary}[theorem]{Corollary}
\theoremstyle{definition}
\newtheorem{definition}[theorem]{Definition}
\newtheorem{remark}[theorem]{Remark}
\newtheorem{helper}[theorem]{Lemma}
\newcommand{\fint}{\mathop{\ooalign{$\int$\cr$-$}}}
\providecommand{\R}{\mathbb{R}}
\providecommand{\N}{\mathbb{N}}
\providecommand{\X}{\mathfrak{X}}
\providecommand{\XX}{\X\X}
\providecommand{\LL}{\mathcal{L}}
\providecommand{\E}{\mathcal{E}}
\providecommand{\Om}{\Omega}
\DeclareMathOperator{\Tr}{Tr}
\DeclareMathOperator*{\argmax}{argmax}
\providecommand{\Penv}{\mathcal{P}}
\providecommand{\Menv}{\mathcal{M}}
\providecommand{\Cenv}{\mathcal{C}}
\providecommand{\Senv}{\mathcal{S}}
\providecommand{\Aclass}{\mathcal{A}}
\DeclareMathOperator{\Lip}{Lip}
\DeclareMathOperator{\osc}{osc}
\providecommand{\noderef}[1]{\hyperref[node:#1]{\texttt{\detokenize{#1}}}}



\title[Blueprint of the Lean formalization]{A note on the comparison principle for degenerate sub-elliptic equations: blueprint of the Lean formalization}
\author{Juan J. Manfredi}
\author{Shirsho Mukherjee}

\begin{document}
\begin{abstract}
This document collects the natural-language statements and detailed proofs that accompany the Lean~4 formalization of Theorem~1.3 and Corollary~1.4 of J.~J.~Manfredi and S.~Mukherjee, \emph{A note on the comparison principle for degenerate sub-elliptic equations}.
The results are formalized for an arbitrary family of smooth vector fields on $\mathbb{R}^n$ satisfying H\"ormander's rank condition; the paper's setting, a basis of the first layer of a Carnot group, is a special case.
The mathematics is that of Section~3.1 of the paper; this document only reorganizes it into the units of the formalization.
Each subsection corresponds to the Lean declaration of the same name in the directory \texttt{C2CompND/}; cross-references point to the results that a proof uses.
Section~\ref{sec:main} lists the formalized results of the paper.
Within each section, items appear in dependency order.
\end{abstract}
\maketitle
\tableofcontents

\section{Definitions}
\subsection{\texorpdfstring{\texttt{Preamble}}{Preamble}}\label{node:Preamble}
\begin{definition}[Ambient space, functions, and derivative conventions]
Throughout, $n,m\in\mathbb{N}$ and the ambient space is $\mathbb{R}^n$ equipped with the standard
Euclidean inner product $\langle\cdot,\cdot\rangle$; vectors of $\mathbb{R}^m$ are written as
functions $\{1,\dots,m\}\to\mathbb{R}$ and $\mathbb{R}^{m\times m}$ denotes the real
$m\times m$ matrices.

Scalar functions and vector fields are taken to be defined on all of $\mathbb{R}^n$:
a function of the paper's class $C^2(\Omega)\cap C(\overline\Omega)$ is represented by a function
$u:\mathbb{R}^n\to\mathbb{R}$ together with the hypotheses that $u$ is of class $C^2$ on $\Omega$
and continuous on $\overline\Omega$. Since every hypothesis and every conclusion below is
restricted to $\Omega$ or to $\overline\Omega$, and since $\Omega$ is open, the values of $u$
outside $\overline\Omega$ never enter, and the derivatives of $u$ at points of $\Omega$ are the
intrinsic ones.

For $u:\mathbb{R}^n\to\mathbb{R}$ and $x\in\mathbb{R}^n$, $Du(x)$ denotes the Fr\'echet derivative
of $u$ at $x$, a linear functional on $\mathbb{R}^n$, identified with the Euclidean gradient
through $Du(x)[z]=\langle\nabla u(x),z\rangle$. By convention $Du(x):=0$ when $u$ is not
differentiable at $x$; the same convention is used for all the derivatives below. The second
derivative is the bilinear form
$$D^2u(x)[z_1,z_2]:=\Big(D\big(y\mapsto Du(y)\big)(x)[z_1]\Big)[z_2],$$
so that, in the matrix notation $D^2u=(\partial_{x_i}\partial_{x_j}u)_{i,j}$ of Section~2.1 of
the paper,
$$D^2u(x)[z_1,z_2]=\langle D^2u(x)\,z_1,z_2\rangle=\sum_{i,j=1}^n
\partial_{x_i}\partial_{x_j}u(x)\,(z_1)_i\,(z_2)_j.$$
For a vector field
$V:\mathbb{R}^n\to\mathbb{R}^n$, $DV(x)$ denotes its Jacobian, so that $DV(x)z$ is the derivative
of $V$ at $x$ in the direction $z$.

We say that $f$ is \emph{of class $C^k$ in a neighbourhood of $x$} if there is an open set
containing $x$ on which $f$ is $k$ times continuously differentiable, and \emph{smooth} means of
class $C^k$ for every $k$. We write $g'(t)$ and $g''(t)$ for the first and second derivatives of a
function $g:\mathbb{R}\to\mathbb{R}$, $\overline{S}$ for the closure of a set $S\subset\mathbb{R}^n$,
$\partial S$ for its topological frontier, and $A^T$ for the transpose of a matrix $A$.
\end{definition}

\begin{definition}[Order on symmetric matrices and the Lie bracket]
For symmetric $A,B\in\mathbb{R}^{m\times m}$ we write $A\leq B$ when $B-A$ is positive
semi-definite, i.e. $\langle(B-A)\xi,\xi\rangle\geq0$ for every $\xi\in\mathbb{R}^m$; this is the
order of Section~2.1 of the paper. A matrix is called \emph{positive semi-definite} only if it is
in addition symmetric, so that "$A$ is positive semi-definite'' always abbreviates "$A$ is
symmetric and $\langle A\xi,\xi\rangle\geq0$ for every $\xi$''.

For vector fields $V,W:\mathbb{R}^n\to\mathbb{R}^n$ the Lie bracket is the vector field
$$[V,W](x):=DW(x)\,V(x)-DV(x)\,W(x).$$
\end{definition}

\begin{proposition}[Weierstrass]
Let $X$ be a topological space, let $K\subset X$ be non-empty and compact and let $f:X\to\mathbb{R}$
be continuous on $K$. Then $f$ attains a maximum over $K$: there is $x_0\in K$ with
$f(y)\leq f(x_0)$ for every $y\in K$. The same holds for the minimum.
\end{proposition}

\begin{proposition}[Closure of a bounded set, and the frontier decomposition]
A bounded subset $\Omega\subset\mathbb{R}^n$ has compact closure $\overline\Omega$. For every
subset $S\subset\mathbb{R}^n$ one has $\overline{S}=S\cup\partial S$.
\end{proposition}

\begin{proposition}[First-order necessary condition at a local extremum]
Let $f:\mathbb{R}^n\to\mathbb{R}$ and let $x_0$ be a local maximum point of $f$. Then
$Df(x_0)=0$. The same holds at a local minimum point.
\end{proposition}

\begin{proposition}[Symmetry of the second derivative]
Let $f:\mathbb{R}^n\to\mathbb{R}$ be of class $C^2$ in a neighbourhood of $x$. Then the bilinear
form $D^2f(x)$ is symmetric: $D^2f(x)[z_1,z_2]=D^2f(x)[z_2,z_1]$ for all
$z_1,z_2\in\mathbb{R}^n$.
\end{proposition}

\begin{proposition}[Second-derivative test, sufficient direction]
Let $g:\mathbb{R}\to\mathbb{R}$ and $t_0\in\mathbb{R}$ be such that $g$ is continuous at $t_0$,
$g'(t_0)=0$ and $g''(t_0)>0$. Then $t_0$ is a local minimum point of $g$.
\end{proposition}

\begin{proposition}[Positive semi-definite square root]
Let $A\in\mathbb{R}^{m\times m}$ be positive semi-definite. Then there is a positive semi-definite
$P\in\mathbb{R}^{m\times m}$ with $P^2=A$.
\end{proposition}

\begin{proposition}[Elementary properties of the trace]
For $A,B,C\in\mathbb{R}^{m\times m}$ and $c\in\mathbb{R}$ one has
$\operatorname{Tr}(A(B+C))=\operatorname{Tr}(AB)+\operatorname{Tr}(AC)$,
$\operatorname{Tr}(A(cB))=c\operatorname{Tr}(AB)$,
$\operatorname{Tr}(AB)=\sum_{i,k}A_{i,k}B_{k,i}$ and $\operatorname{Tr}(AB)=\operatorname{Tr}(BA)$.
\end{proposition}


\subsection{\texorpdfstring{\texttt{HGrad}}{HGrad}}\label{node:HGrad}
\begin{definition}
Let $\sigma=(\sigma^1,\dots,\sigma^m)$ be vector fields $\sigma^i:\mathbb{R}^n\to\mathbb{R}^n$ and
let $X_i:=\sigma^i(x)\cdot\nabla$ be the associated first-order operators, as in Section~2.1 of the
paper. For $u:\mathbb{R}^n\to\mathbb{R}$ and $x\in\mathbb{R}^n$, the \emph{horizontal
(sub-elliptic) gradient} of $u$ at $x$ is the vector $\mathfrak{X}u(x)\in\mathbb{R}^m$ with
components
$$\big(\mathfrak{X}u(x)\big)_i\ :=\ Du(x)\big[\sigma^i(x)\big]\ =\ \big\langle\nabla
u(x),\sigma^i(x)\big\rangle\ =\ X_iu(x),\qquad i=1,\dots,m,$$
where $Du(x)$ is the Fr\'echet derivative of $u$ at $x$ (equal to $0$ when $u$ is not
differentiable at $x$, by the convention of the preamble).

This is the quantity $\mathfrak{X}u=(X_1u,\dots,X_mu)$ of Section~2.1 of the paper, equivalently
$\mathfrak{X}u=\sigma(x)^T\nabla u$ in the notation of \texttt{eq:sigrad}, written as a function
$\{1,\dots,m\}\to\mathbb{R}$.
\end{definition}


\subsection{\texorpdfstring{\texttt{HHess}}{HHess}}\label{node:HHess}
\begin{definition}
Let $\sigma=(\sigma^1,\dots,\sigma^m)$ be vector fields $\sigma^i:\mathbb{R}^n\to\mathbb{R}^n$,
with associated first-order operators $X_i=\sigma^i(x)\cdot\nabla$, and let
$u:\mathbb{R}^n\to\mathbb{R}$ and $x\in\mathbb{R}^n$. The \emph{horizontal second derivative
matrix} of $u$ at $x$ is the matrix
$\mathfrak{X}\mathfrak{X}u(x)\in\mathbb{R}^{m\times m}$ defined entrywise by
$$\big(\mathfrak{X}\mathfrak{X}u(x)\big)_{i,j}\ :=\ D\Big(y\mapsto
Du(y)\big[\sigma^i(y)\big]\Big)(x)\big[\sigma^j(x)\big]\ =\ X_jX_iu(x),
\qquad i,j=1,\dots,m,$$
that is, the $(i,j)$ entry is obtained by differentiating the function
$y\mapsto X_iu(y)=\big(\mathfrak{X}u(y)\big)_i$ at $x$ in the direction $\sigma^j(x)$. All
derivatives follow the convention of the preamble (they are $0$ where the relevant function fails
to be differentiable).

Thus the $(i,j)$ entry is $X_jX_iu$, matching the convention
$\mathfrak{X}\mathfrak{X}u=(X_jX_iu)_{i,j}$ of Section~2.1 of the paper. This matrix is not
symmetric in general.
\end{definition}


\subsection{\texorpdfstring{\texttt{HHessSym}}{HHessSym}}\label{node:HHessSym}
\begin{definition}
Let $\sigma$, $u$ and $x$ be as in \noderef{HHess}. The \emph{symmetrized horizontal second
derivative} of $u$ at $x$ is the matrix
$$\mathfrak{X}\mathfrak{X}^\star u(x)\ :=\ \tfrac{1}{2}\Big(\mathfrak{X}\mathfrak{X}u(x)+
\big(\mathfrak{X}\mathfrak{X}u(x)\big)^T\Big)\ \in\ \mathbb{R}^{m\times m},$$
whose entries are
$\big(\mathfrak{X}\mathfrak{X}^\star u(x)\big)_{i,j}=\tfrac12\big(X_jX_iu(x)+X_iX_ju(x)\big)$.
By construction $\mathfrak{X}\mathfrak{X}^\star u(x)$ is symmetric.

This is the quantity
$\mathfrak{X}\mathfrak{X}^\star u=\tfrac12\big(\mathfrak{X}\mathfrak{X}u+(\mathfrak{X}\mathfrak{X}
u)^T\big)=\tfrac12\big(X_iX_ju+X_jX_iu\big)_{i,j}$ of Section~2.1 of the paper.
\end{definition}


\subsection{\texorpdfstring{\texttt{SubellipticOp}}{SubellipticOp}}\label{node:SubellipticOp}
\begin{definition}
Let $\sigma=(\sigma^1,\dots,\sigma^m)$ be vector fields $\sigma^i:\mathbb{R}^n\to\mathbb{R}^n$ and
let $a:\mathbb{R}^m\to\mathbb{R}^{m\times m}$, $a(\xi)=(a_{i,j}(\xi))_{i,j}$. For
$u:\mathbb{R}^n\to\mathbb{R}$ and $x\in\mathbb{R}^n$ we set
$$\mathcal{L}u(x)\ :=\ -\operatorname{Tr}\Big(a\big(\mathfrak{X}u(x)\big)\,
\mathfrak{X}\mathfrak{X}u(x)\Big)\ =\ -\sum_{i,j=1}^m
a_{i,j}\big(\mathfrak{X}u(x)\big)\,X_iX_ju(x),$$
where $\mathfrak{X}u(x)$ is the horizontal gradient of \noderef{HGrad} and
$\mathfrak{X}\mathfrak{X}u(x)$ the horizontal second derivative matrix of \noderef{HHess}; the
second expression is the identity $\operatorname{Tr}(AB)=\sum_{i,k}A_{i,k}B_{k,i}$ together with
$\big(\mathfrak{X}\mathfrak{X}u\big)_{j,i}=X_iX_ju$.

This is the operator of \texttt{eq:maineq} of the paper. Whenever $a(\xi)$ is symmetric the second
expression coincides, after exchanging the summation indices, with the paper's
$-\sum_{i,j}a_{i,j}(\mathfrak{X}u)X_jX_iu$.
\end{definition}


\subsection{\texorpdfstring{\texttt{EllipticityFn}}{EllipticityFn}}\label{node:EllipticityFn}
\begin{definition}
Let $a:\mathbb{R}^m\to\mathbb{R}^{m\times m}$, $a(\xi)=(a_{i,j}(\xi))_{i,j}$. The associated
\emph{ellipticity function} $\mathcal{E}:\mathbb{R}^m\to\mathbb{R}$ is
$$\mathcal{E}(\xi)\ :=\ \big\langle a(\xi)\,\xi,\xi\big\rangle\ =\ \sum_{i,j=1}^m
\xi_i\,a_{i,j}(\xi)\,\xi_j .$$
This is the quantity $\mathcal{E}(\xi)=\langle a(\xi)\xi,\xi\rangle$ appearing in the ellipticity
condition \texttt{eq:ell} of the paper.
\end{definition}


\subsection{\texorpdfstring{\texttt{IsClassicalSubsol}}{IsClassicalSubsol}}\label{node:IsClassicalSubsol}
\begin{definition}
Let $\sigma$ and $a$ be as in \noderef{SubellipticOp}, let $U\subset\mathbb{R}^n$ and let
$u:\mathbb{R}^n\to\mathbb{R}$. We say that $u$ is a \emph{classical sub-solution} of
\texttt{eq:maineq} in $U$, and write $\mathcal{L}u\leq0$ in $U$, if
$$\mathcal{L}u(x)\leq0\qquad\text{at every }x\in U\text{ with }\mathfrak{X}u(x)\neq0 .$$
Here $\mathfrak{X}u$ is the horizontal gradient of \noderef{HGrad} and $\mathcal{L}$ the operator
of \noderef{SubellipticOp}. No condition is imposed at the horizontal critical points of $u$,
where $a$ need not be defined.

This is the sub-solution condition of Definition \texttt{def:classical} of the paper. The paper
additionally takes $U$ open and $u\in C^2(U)$; those are regularity requirements on $U$ and $u$
rather than part of the differential inequality, and the nodes that use this notion carry them as
separate hypotheses.
\end{definition}


\subsection{\texorpdfstring{\texttt{IsClassicalSupersol}}{IsClassicalSupersol}}\label{node:IsClassicalSupersol}
\begin{definition}
Let $\sigma$ and $a$ be as in \noderef{SubellipticOp}, let $U\subset\mathbb{R}^n$ and let
$u:\mathbb{R}^n\to\mathbb{R}$. We say that $u$ is a \emph{classical super-solution} of
\texttt{eq:maineq} in $U$, and write $\mathcal{L}u\geq0$ in $U$, if
$$\mathcal{L}u(x)\geq0\qquad\text{at every }x\in U\text{ with }\mathfrak{X}u(x)\neq0 .$$
Here $\mathfrak{X}u$ is the horizontal gradient of \noderef{HGrad} and $\mathcal{L}$ the operator
of \noderef{SubellipticOp}. No condition is imposed at the horizontal critical points of $u$,
where $a$ need not be defined.

This is the super-solution condition of Definition \texttt{def:classical} of the paper, which
defines sub- and super-solutions in one breath. It is not a restatement of the sub-solution
condition of \noderef{IsClassicalSubsol} applied to $-u$: the matrix $a$ is not assumed to be even
(Remark \texttt{rem:hom}(a) of the paper says so explicitly), so in general
$a(-\mathfrak{X}u)\neq a(\mathfrak{X}u)$ and $\mathcal{L}(-u)\neq-\mathcal{L}u$. As in
\noderef{IsClassicalSubsol}, the paper's openness of $U$ and $C^2$ regularity of $u$ are carried as
separate hypotheses by the nodes that use this notion.
\end{definition}


\subsection{\texorpdfstring{\texttt{HorFamily}}{HorFamily}}\label{node:HorFamily}
\begin{definition}
Let $n,m\in\mathbb{N}$ and let $\sigma=(\sigma^1,\dots,\sigma^m)$ be a family of vector fields
$\sigma^i:\mathbb{R}^n\to\mathbb{R}^n$. We define $\mathrm{HorFamily}(\sigma)$, a family of vector
fields on $\mathbb{R}^n$, inductively by the two generation rules:
\begin{itemize}
\item[(i)] $\sigma^i\in\mathrm{HorFamily}(\sigma)$ for every $i\in\{1,\dots,m\}$;
\item[(ii)] if $V,W\in\mathrm{HorFamily}(\sigma)$, then $[V,W]\in\mathrm{HorFamily}(\sigma)$,
where $[V,W](x)=DW(x)\,V(x)-DV(x)\,W(x)$ is the Lie bracket.
\end{itemize}
Equivalently, $\mathrm{HorFamily}(\sigma)$ is the smallest family of vector fields on
$\mathbb{R}^n$ that contains each $\sigma^i$ and is closed under the Lie bracket.

This is the family of iterated commutators denoted $\mathrm{Lie}(X_1,\dots,X_m)$ in the paper's
H\"ormander rank condition \texttt{eq:horcond}, written for the first-order operators
$X_i=\sigma^i(x)\cdot\nabla$ in terms of the coefficient fields $\sigma^i$ of Section~2.1 of the
paper.
\end{definition}


\subsection{\texorpdfstring{\texttt{HormanderRank}}{HormanderRank}}\label{node:HormanderRank}
\begin{definition}
Let $\sigma=(\sigma^1,\dots,\sigma^m)$ be a family of vector fields
$\sigma^i:\mathbb{R}^n\to\mathbb{R}^n$. We say that $\sigma$ \emph{satisfies the H\"ormander rank
condition}, written $\mathrm{HormanderRank}(\sigma)$, if for every $x\in\mathbb{R}^n$ the linear
span of the set of evaluations at $x$ of the vector fields of $\mathrm{HorFamily}(\sigma)$ (see
\noderef{HorFamily}) is the whole of $\mathbb{R}^n$:
$$\mathrm{span}_{\mathbb{R}}\big\{V(x)\ :\ V\in\mathrm{HorFamily}(\sigma)\big\}=\mathbb{R}^n
\qquad\text{for every }x\in\mathbb{R}^n .$$
Since $\mathbb{R}^n$ has dimension $n$, this says exactly that
$\dim\big(\mathrm{Lie}(X_1,\dots,X_m)(x)\big)=n$ at every point $x\in\mathbb{R}^n$, which is the
paper's condition \texttt{eq:horcond} for the first-order operators
$X_i=\sigma^i(x)\cdot\nabla$.
\end{definition}


\section{Lemmas}
\subsection{\texorpdfstring{\texttt{FDerivEqAtLocalMaxSub}}{FDerivEqAtLocalMaxSub}}\label{node:FDerivEqAtLocalMaxSub}
\begin{lemma}
Let $u,v:\mathbb{R}^n\to\mathbb{R}$ and let $x_0\in\mathbb{R}^n$. Assume that $u$ is differentiable
at $x_0$, that $v$ is differentiable at $x_0$, and that $x_0$ is a local maximum point of $u-v$.
Then
$$Du(x_0)=Dv(x_0)$$
as linear functionals on $\mathbb{R}^n$. This is the first clause of the sentence opening the proof
of Lemma \texttt{lem:vweak} of the paper, namely that the interior maximality at $x_0$ implies
$\nabla u(x_0)=\nabla v(x_0)$; the paper derives the horizontal identity
$\mathfrak{X}u(x_0)=\mathfrak{X}v(x_0)$ from it through \texttt{eq:sigrad}.
\end{lemma}
\begin{proof}
Since $u$ and $v$ are differentiable at $x_0$, so is $u-v$, and the derivative of a difference is
the difference of the derivatives:
$$D(u-v)(x_0)=Du(x_0)-Dv(x_0).$$
As $x_0$ is a local maximum point of $u-v$, the first-order necessary condition at a local extremum
(preamble) applied to $f:=u-v$ gives
$$D(u-v)(x_0)=0.$$
Combining the two displays, $Du(x_0)-Dv(x_0)=0$ in the vector space of linear functionals on
$\mathbb{R}^n$, and therefore $Du(x_0)=Dv(x_0)$.
\end{proof}


\subsection{\texorpdfstring{\texttt{HGradEqAtLocalMax}}{HGradEqAtLocalMax}}\label{node:HGradEqAtLocalMax}
\begin{lemma}
Let $\sigma=(\sigma^1,\dots,\sigma^m)$ be vector fields $\sigma^i:\mathbb{R}^n\to\mathbb{R}^n$, let
$u,v:\mathbb{R}^n\to\mathbb{R}$ and let $x_0\in\mathbb{R}^n$. Assume that $u$ and $v$ are
differentiable at $x_0$ and that $x_0$ is a local maximum point of $u-v$. Then
$$\mathfrak{X}u(x_0)=\mathfrak{X}v(x_0).$$
\end{lemma}
\begin{proof}
Since $u$ and $v$ are differentiable at $x_0$ and $x_0$ is a local maximum point of $u-v$, the
hypotheses of \noderef{FDerivEqAtLocalMaxSub} are met, and it gives
$$Du(x_0)=Dv(x_0)$$
as linear functionals on $\mathbb{R}^n$.

Evaluating this equality of functionals at the vector $\sigma^i(x_0)\in\mathbb{R}^n$ gives, for
every $i\in\{1,\dots,m\}$,
$$\big(\mathfrak{X}u(x_0)\big)_i=Du(x_0)\big[\sigma^i(x_0)\big]
=Dv(x_0)\big[\sigma^i(x_0)\big]=\big(\mathfrak{X}v(x_0)\big)_i,$$
by the definition of the horizontal gradient (\noderef{HGrad}). Hence
$\mathfrak{X}u(x_0)=\mathfrak{X}v(x_0)$.
\end{proof}


\subsection{\texorpdfstring{\texttt{DerivDerivRestrictLine}}{DerivDerivRestrictLine}}\label{node:DerivDerivRestrictLine}
\begin{lemma}
Let $n\in\mathbb{N}$, let $f:\mathbb{R}^n\to\mathbb{R}$, let $x_0\in\mathbb{R}^n$ and let
$w\in\mathbb{R}^n$. Assume that $f$ is of class $C^2$ in a neighbourhood of $x_0$, and let
$g:\mathbb{R}\to\mathbb{R}$, $g(t):=f(x_0+tw)$, be the restriction of $f$ to the line through
$x_0$ in the direction $w$. Then
$$g''(0)=D^2f(x_0)[w,w].$$
\end{lemma}
\begin{proof}
Let $\gamma(t):=x_0+tw$, so that $g=f\circ\gamma$. The map $\gamma:\mathbb{R}\to\mathbb{R}^n$ is
affine, hence smooth, with $\gamma(0)=x_0$ and $\gamma'(t)=w$ for every $t$; being continuous with
$\gamma(0)=x_0$, it satisfies $\gamma(t)\to x_0$ as $t\to0$.

\emph{The first derivative of $g$ near $0$.} By hypothesis there is an open set $V\ni x_0$ on which
$f$ is of class $C^2$; in particular $f$ is differentiable at every point of $V$. Since
$\gamma(t)\to x_0$ as $t\to0$, there is a neighbourhood $N$ of $0$ with $\gamma(t)\in V$ for every
$t\in N$. For such $t$ the chain rule applies to $f\circ\gamma$ at $t$ and gives
\begin{equation*}g'(t)=Df(\gamma(t))\big[\gamma'(t)\big]=Df(x_0+tw)[w]\qquad\text{for every }t\in N. \tag{1}\end{equation*}

\emph{Differentiating once more at $0$.} Since $f$ is of class $C^2$ on $V$, the map
$y\mapsto Df(y)$ is differentiable at $x_0$, with derivative the linear map
$z\mapsto D\big(y\mapsto Df(y)\big)(x_0)[z]$. Composing it with the affine map $\gamma$ (chain
rule again, using $\gamma(0)=x_0$ and $\gamma'(0)=w$), the map $t\mapsto Df(x_0+tw)$ is
differentiable at $t=0$ with derivative $D\big(y\mapsto Df(y)\big)(x_0)[w]$. Composing further
with the evaluation map $\Lambda\mapsto\Lambda[w]$, which is linear and hence smooth, the
real-valued function $t\mapsto Df(x_0+tw)[w]$ is differentiable at $t=0$ with derivative
$$\Big(D\big(y\mapsto Df(y)\big)(x_0)[w]\Big)[w]=D^2f(x_0)[w,w],$$
the last equality being the definition of $D^2f(x_0)$ fixed in the preamble.

\emph{Conclusion.} By (1) the functions $g'$ and $t\mapsto Df(x_0+tw)[w]$ agree on the
neighbourhood $N$ of $0$, so they have the same derivative at $0$. Hence $g'$ is differentiable at
$0$ and $g''(0)=D^2f(x_0)[w,w]$.
\end{proof}


\subsection{\texorpdfstring{\texttt{SecondDerivTestOneDim}}{SecondDerivTestOneDim}}\label{node:SecondDerivTestOneDim}
\begin{lemma}
Let $g:\mathbb{R}\to\mathbb{R}$ and $t_0\in\mathbb{R}$. Assume that $g$ is of class $C^2$ in a
neighbourhood of $t_0$ and that $t_0$ is a local maximum point of $g$. Then
$$g''(t_0)\leq0 .$$
\end{lemma}
\begin{proof}
Suppose, for contradiction, that $g''(t_0)>0$.

Since $g$ is of class $C^2$ near $t_0$, it is in particular differentiable at $t_0$, and $t_0$ is a
local maximum point of $g$, so the first-order necessary condition (preamble) gives $g'(t_0)=0$.
Since $g$ is continuous at $t_0$, $g'(t_0)=0$ and $g''(t_0)>0$, the sufficient direction of the
second-derivative test (preamble) shows that $t_0$ is a \emph{local minimum} point of $g$.

Thus $t_0$ is both a local maximum and a local minimum point of $g$: there are neighbourhoods
$N_1,N_2$ of $t_0$ with $g\leq g(t_0)$ on $N_1$ and $g\geq g(t_0)$ on $N_2$. On the neighbourhood
$N:=N_1\cap N_2$ of $t_0$ we therefore have $g\equiv g(t_0)$. Let $I\subset N$ be an open interval
containing $t_0$. Then $g$ is constant on $I$, hence $g'\equiv0$ on $I$, i.e.\ $g'$ coincides with
the zero function on a neighbourhood of $t_0$. Differentiating, $g''(t_0)=(g')'(t_0)=0$, which
contradicts $g''(t_0)>0$.

Hence $g''(t_0)\leq0$.
\end{proof}


\subsection{\texorpdfstring{\texttt{SndFDerivNonposAtLocalMax}}{SndFDerivNonposAtLocalMax}}\label{node:SndFDerivNonposAtLocalMax}
\begin{lemma}
Let $n\in\mathbb{N}$, let $f:\mathbb{R}^n\to\mathbb{R}$ and let $x_0\in\mathbb{R}^n$. Assume that
$f$ is of class $C^2$ in a neighbourhood of $x_0$ and that $x_0$ is a local maximum point of $f$.
Then
$$D^2f(x_0)[w,w]\ \leq\ 0\qquad\text{for every }w\in\mathbb{R}^n .$$
\end{lemma}
\begin{proof}
Fix $w\in\mathbb{R}^n$ and let $\gamma(t):=x_0+tw$, $g:=f\circ\gamma$, i.e.\ $g(t)=f(x_0+tw)$ for
$t\in\mathbb{R}$. The map $\gamma:\mathbb{R}\to\mathbb{R}^n$ is affine, hence smooth, with
$\gamma(0)=x_0$ and $\gamma'(t)=w$ for every $t$.

\emph{Step 1: $g$ is of class $C^2$ in a neighbourhood of $0$.} By hypothesis there is an open set
$V\ni x_0$ on which $f$ is of class $C^2$. Since $\gamma$ is continuous and $\gamma(0)=x_0$, the
set $\gamma^{-1}(V)$ is open and contains $0$; on it $g=f\circ\gamma$ is the composition of a
smooth map with a $C^2$ map, hence of class $C^2$.

\emph{Step 2: $0$ is a local maximum point of $g$.} By hypothesis there is a neighbourhood
$N\ni x_0$ with $f(y)\leq f(x_0)$ for all $y\in N$. Then $\gamma^{-1}(N)$ is a neighbourhood of
$0$, and for $t\in\gamma^{-1}(N)$ we get $g(t)=f(\gamma(t))\leq f(x_0)=g(0)$.

\emph{Step 3: $g''(0)=D^2f(x_0)[w,w]$.} The function $f$ is of class $C^2$ in a neighbourhood of
$x_0$, so \noderef{DerivDerivRestrictLine} applies to $f$, $x_0$ and $w$. Since $g$ is precisely
the restriction $t\mapsto f(x_0+tw)$ of $f$ to the line through $x_0$ in the direction $w$, its
conclusion reads $g''(0)=D^2f(x_0)[w,w]$.

\emph{Conclusion.} By Steps 1 and 2 and \noderef{SecondDerivTestOneDim} applied to $g$ at
$t_0=0$, we have $g''(0)\leq0$. By Step 3 this is exactly $D^2f(x_0)[w,w]\leq0$.
\end{proof}


\subsection{\texorpdfstring{\texttt{SndFDerivLeAtLocalMaxSub}}{SndFDerivLeAtLocalMaxSub}}\label{node:SndFDerivLeAtLocalMaxSub}
\begin{lemma}
Let $U\subset\mathbb{R}^n$ be open, let $u,v:\mathbb{R}^n\to\mathbb{R}$ be of class $C^2$ on $U$
and let $x_0\in U$ be a local maximum point of $u-v$. Then
$$D^2u(x_0)[w,w]\ \leq\ D^2v(x_0)[w,w]\qquad\text{for every }w\in\mathbb{R}^n,$$
i.e.\ $D^2u(x_0)\leq D^2v(x_0)$. This is the second-order clause of the sentence opening the proof
of Lemma \texttt{lem:vweak} of the paper, that the interior maximality at $x_0$ implies
$\nabla u(x_0)=\nabla v(x_0)$ and $D^2u(x_0)\leq D^2v(x_0)$; the first-order clause of that same
sentence is the node \texttt{FDerivEqAtLocalMaxSub}.
\end{lemma}
\begin{proof}
Put $f:=u-v$. Since $U$ is open, $x_0\in U$ and $u,v$ are of class $C^2$ on $U$, both $u$ and $v$
are of class $C^2$ in a neighbourhood of $x_0$, namely on $U$ itself; hence so is $f$, and in
particular $u$, $v$ and $f$ are differentiable at every point of $U$.

\emph{Step 1: the first derivatives subtract on $U$.} For every $y\in U$ both $u$ and $v$ are
differentiable at $y$, so $f=u-v$ is differentiable at $y$ with
\begin{equation*}Df(y)=Du(y)-Dv(y)\qquad\text{for every }y\in U. \tag{1}\end{equation*}
Since $U$ is a neighbourhood of $x_0$, the maps $y\mapsto Df(y)$ and $y\mapsto Du(y)-Dv(y)$
therefore coincide on a neighbourhood of $x_0$.

\emph{Step 2: the second derivatives subtract at $x_0$.} Because $u$ and $v$ are of class $C^2$ in
a neighbourhood of $x_0$, the maps $y\mapsto Du(y)$ and $y\mapsto Dv(y)$ are differentiable at
$x_0$, hence so is their difference. By Step 1 the map $y\mapsto Df(y)$ agrees with that difference
on a neighbourhood of $x_0$, so it has the same derivative at $x_0$, and by linearity of the
derivative
\begin{equation*}D^2f(x_0)[z_1,z_2]=D^2u(x_0)[z_1,z_2]-D^2v(x_0)[z_1,z_2]\qquad
\text{for all }z_1,z_2\in\mathbb{R}^n. \tag{2}\end{equation*}

\emph{Step 3: conclusion.} The function $f$ is of class $C^2$ in a neighbourhood of $x_0$ and has a
local maximum at $x_0$, so \noderef{SndFDerivNonposAtLocalMax} applies to $f$ and gives
$$D^2f(x_0)[w,w]\leq0\qquad\text{for every }w\in\mathbb{R}^n .$$
Fix $w\in\mathbb{R}^n$ and combine this with (2) taken at $z_1=z_2=w$:
$$D^2u(x_0)[w,w]-D^2v(x_0)[w,w]=D^2f(x_0)[w,w]\leq0,$$
that is, $D^2u(x_0)[w,w]\leq D^2v(x_0)[w,w]$.
\end{proof}


\subsection{\texorpdfstring{\texttt{HHessEntryFormula}}{HHessEntryFormula}}\label{node:HHessEntryFormula}
\begin{lemma}
Let $\sigma=(\sigma^1,\dots,\sigma^m)$ be vector fields $\sigma^i:\mathbb{R}^n\to\mathbb{R}^n$, let
$u:\mathbb{R}^n\to\mathbb{R}$, let $x\in\mathbb{R}^n$ and let $i,j\in\{1,\dots,m\}$. Assume that
$u$ is of class $C^2$ in a neighbourhood of $x$ and that $\sigma^i$ is of class $C^1$ in a
neighbourhood of $x$. Then
$$\big(\mathfrak{X}\mathfrak{X}u(x)\big)_{i,j}\ =\
D^2u(x)\big[\sigma^j(x),\sigma^i(x)\big]\ +\ Du(x)\big[D\sigma^i(x)\,\sigma^j(x)\big],$$
that is, $X_jX_iu=\langle D^2u\,\sigma^j,\sigma^i\rangle+\langle\nabla
u,D\sigma^i\,\sigma^j\rangle$. This is the identity $X_jX_iu=\langle
D^2u\,\sigma^i,\sigma^j\rangle+\langle D\sigma^i\,\sigma^j,\nabla u\rangle$ recorded in
Section~2.1 of the paper, which underlies \texttt{eq:sigrad} and the formula \texttt{eq:Mxz} for
the first-order matrix $\mathcal{M}(x,\zeta)$.
\end{lemma}
\begin{proof}
By the definition of the horizontal second derivative matrix (\noderef{HHess}),
$$\big(\mathfrak{X}\mathfrak{X}u(x)\big)_{i,j}=D\varphi(x)\big[\sigma^j(x)\big],\qquad\text{where}
\qquad \varphi(y):=Du(y)\big[\sigma^i(y)\big].$$
It therefore suffices to compute $D\varphi(x)[z]$ for an arbitrary $z\in\mathbb{R}^n$ and then take
$z=\sigma^j(x)$.

The function $\varphi$ is the composition of the map
$y\mapsto\big(Du(y),\sigma^i(y)\big)$ with the evaluation pairing
$$\Phi:\big(\mathbb{R}^n\big)^*\times\mathbb{R}^n\to\mathbb{R},\qquad
\Phi(\Lambda,\eta):=\Lambda[\eta],$$
which is bilinear and continuous. Its two components are differentiable at $x$:
\begin{itemize}
\item since $u$ is of class $C^2$ in a neighbourhood of $x$, the map $y\mapsto Du(y)$ is
differentiable at $x$, with derivative $z\mapsto D\big(y\mapsto Du(y)\big)(x)[z]$;
\item since $\sigma^i$ is of class $C^1$ in a neighbourhood of $x$, the map $y\mapsto\sigma^i(y)$
is differentiable at $x$, with derivative $z\mapsto D\sigma^i(x)\,z$.
\end{itemize}
By the chain rule applied to $\Phi$ composed with this pair (equivalently, by the product rule for
a continuous bilinear pairing), $\varphi$ is differentiable at $x$ and for every
$z\in\mathbb{R}^n$
$$D\varphi(x)[z]=\Big(D\big(y\mapsto Du(y)\big)(x)[z]\Big)\big[\sigma^i(x)\big]
\ +\ Du(x)\big[D\sigma^i(x)\,z\big].$$
By the definition of the second derivative fixed in the preamble, the first summand is
$D^2u(x)\big[z,\sigma^i(x)\big]$. Hence
$$D\varphi(x)[z]=D^2u(x)\big[z,\sigma^i(x)\big]+Du(x)\big[D\sigma^i(x)\,z\big].$$
Taking $z=\sigma^j(x)$ gives exactly the asserted identity.
\end{proof}


\subsection{\texorpdfstring{\texttt{PosSemidefSymmQuadForm}}{PosSemidefSymmQuadForm}}\label{node:PosSemidefSymmQuadForm}
\begin{lemma}
Let $n,m\in\mathbb{N}$, let $\mathcal{W}$ be a continuous bilinear form on $\mathbb{R}^n$ such that
$$\mathcal{W}[z,z]\ \geq\ 0\qquad\text{for every }z\in\mathbb{R}^n,$$
and let $s^1,\dots,s^m\in\mathbb{R}^n$ be arbitrary. Then the matrix
$P\in\mathbb{R}^{m\times m}$ with entries
$$P_{i,j}:=\tfrac12\Big(\mathcal{W}\big[s^j,s^i\big]+\mathcal{W}\big[s^i,s^j\big]\Big)$$
is positive semi-definite, i.e.\ symmetric with $\langle P\xi,\xi\rangle\geq0$ for every
$\xi\in\mathbb{R}^m$.
\end{lemma}
\begin{proof}
\emph{Step 1: $P$ is symmetric.} Exchanging $i$ and $j$ in the defining formula exchanges the two
summands and leaves their sum unchanged:
$$P_{j,i}=\tfrac12\Big(\mathcal{W}\big[s^i,s^j\big]+\mathcal{W}\big[s^j,s^i\big]\Big)=P_{i,j}.$$

\emph{Step 2: expansion of the quadratic form.} Fix $\xi\in\mathbb{R}^m$ and put
$z:=\sum_{i=1}^m\xi_i\,s^i\in\mathbb{R}^n$. By the definition of $P$ and bilinearity of the
summation,
\begin{equation*}\langle P\xi,\xi\rangle=\sum_{i,j=1}^m\xi_i\,P_{i,j}\,\xi_j
=\tfrac12\sum_{i,j=1}^m\xi_i\xi_j\,\mathcal{W}\big[s^j,s^i\big]
+\tfrac12\sum_{i,j=1}^m\xi_i\xi_j\,\mathcal{W}\big[s^i,s^j\big]. \tag{1}\end{equation*}

\emph{Step 3: the two double sums agree.} Renaming the pair of summation indices
$i\leftrightarrow j$ in the first double sum of (1), and using $\xi_j\xi_i=\xi_i\xi_j$,
\begin{equation*}\sum_{i,j=1}^m\xi_i\xi_j\,\mathcal{W}\big[s^j,s^i\big]
=\sum_{i,j=1}^m\xi_j\xi_i\,\mathcal{W}\big[s^i,s^j\big]
=\sum_{i,j=1}^m\xi_i\xi_j\,\mathcal{W}\big[s^i,s^j\big]. \tag{2}\end{equation*}
(Both double sums are finite, so the renaming is a bijective reindexing of the index set
$\{1,\dots,m\}^2$.) Substituting (2) into (1), the two halves coincide and
$$\langle P\xi,\xi\rangle=\sum_{i,j=1}^m\xi_i\xi_j\,\mathcal{W}\big[s^j,s^i\big].$$

\emph{Step 4: conclusion.} By bilinearity of $\mathcal{W}$, expanding $z$ in both arguments,
$$\mathcal{W}[z,z]=\sum_{i,j=1}^m\xi_i\xi_j\,\mathcal{W}\big[s^i,s^j\big]
=\sum_{i,j=1}^m\xi_i\xi_j\,\mathcal{W}\big[s^j,s^i\big],$$
the second equality being (2) read backwards. Hence
$\langle P\xi,\xi\rangle=\mathcal{W}[z,z]\geq0$ by hypothesis. Since $\xi\in\mathbb{R}^m$ was
arbitrary, this together with the symmetry of Step 1 says that $P$ is positive semi-definite.
\end{proof}


\subsection{\texorpdfstring{\texttt{HGradHHessComparison}}{HGradHHessComparison}}\label{node:HGradHHessComparison}
\begin{lemma}
Let $\sigma=(\sigma^1,\dots,\sigma^m)$ be vector fields $\sigma^i:\mathbb{R}^n\to\mathbb{R}^n$, let
$u,v:\mathbb{R}^n\to\mathbb{R}$ and let $x\in\mathbb{R}^n$. Assume that $u$ and $v$ are of class
$C^2$ in a neighbourhood of $x$, that every $\sigma^i$ is of class $C^1$ in a neighbourhood of
$x$, and that
$$Du(x)=Dv(x)\qquad\text{and}\qquad D^2u(x)[w,w]\leq D^2v(x)[w,w]\quad\text{for every }
w\in\mathbb{R}^n .$$
Then
$$\mathfrak{X}u(x)=\mathfrak{X}v(x)\qquad\text{and}\qquad
\mathfrak{X}\mathfrak{X}^\star v(x)-\mathfrak{X}\mathfrak{X}^\star u(x)\ \text{ is positive
semi-definite},$$
the latter being $\mathfrak{X}\mathfrak{X}^\star u(x)\leq\mathfrak{X}\mathfrak{X}^\star v(x)$ in the
order of Section~2.1 of the paper. This is the last sentence of Section~2.1 of the paper: if
$u,w\in C^2$ satisfy $\nabla u(x)=\nabla w(x)$ and $D^2u(x)\leq D^2w(x)$ at a point $x$, then
$\mathfrak{X}u(x)=\mathfrak{X}w(x)$ and
$\mathfrak{X}\mathfrak{X}^\star u(x)\leq\mathfrak{X}\mathfrak{X}^\star w(x)$.
\end{lemma}
\begin{proof}
\emph{Step 1: the horizontal gradients.} Evaluating the equality of linear functionals
$Du(x)=Dv(x)$ at $\sigma^i(x)$ gives, for every $i$,
$\big(\mathfrak{X}u(x)\big)_i=Du(x)[\sigma^i(x)]=Dv(x)[\sigma^i(x)]=\big(\mathfrak{X}v(x)\big)_i$
by the definition of \noderef{HGrad}. Hence $\mathfrak{X}u(x)=\mathfrak{X}v(x)$.

\emph{Step 2: the entries of the horizontal second derivative matrices.} Fix $i,j$. The hypotheses
of \noderef{HHessEntryFormula} hold for $u$ (of class $C^2$ near $x$) and for $v$ (of class $C^2$
near $x$), with $\sigma^i$ of class $C^1$ near $x$. Therefore
$$\big(\mathfrak{X}\mathfrak{X}u(x)\big)_{i,j}
=D^2u(x)\big[\sigma^j(x),\sigma^i(x)\big]+Du(x)\big[D\sigma^i(x)\,\sigma^j(x)\big],$$
$$\big(\mathfrak{X}\mathfrak{X}v(x)\big)_{i,j}
=D^2v(x)\big[\sigma^j(x),\sigma^i(x)\big]+Dv(x)\big[D\sigma^i(x)\,\sigma^j(x)\big].$$
Since $Du(x)=Dv(x)$, the two first-order terms are equal, and subtracting gives
\begin{equation*}\big(\mathfrak{X}\mathfrak{X}v(x)\big)_{i,j}-\big(\mathfrak{X}\mathfrak{X}u(x)\big)_{i,j}
=\mathcal{W}\big[\sigma^j(x),\sigma^i(x)\big],\qquad
\mathcal{W}:=D^2v(x)-D^2u(x), \tag{1}\end{equation*}
where $\mathcal{W}$ is a bilinear form on $\mathbb{R}^n$.

\emph{Step 3: the difference of the symmetrized matrices.} Write
$P:=\mathfrak{X}\mathfrak{X}^\star v(x)-\mathfrak{X}\mathfrak{X}^\star u(x)$ and
$\Delta:=\mathfrak{X}\mathfrak{X}v(x)-\mathfrak{X}\mathfrak{X}u(x)$. By the definition of the
symmetrization (\noderef{HHessSym}) and linearity, $P=\tfrac12\big(\Delta+\Delta^T\big)$, so that
by (1)
\begin{equation*}P_{i,j}=\tfrac12\big(\Delta_{i,j}+\Delta_{j,i}\big)
=\tfrac12\Big(\mathcal{W}\big[\sigma^j(x),\sigma^i(x)\big]
+\mathcal{W}\big[\sigma^i(x),\sigma^j(x)\big]\Big) \tag{2}\end{equation*}
for all $i,j\in\{1,\dots,m\}$.

\emph{Step 4: $P$ is positive semi-definite.} The form $\mathcal{W}=D^2v(x)-D^2u(x)$ is a
continuous bilinear form on $\mathbb{R}^n$, being the difference of two second Fr\'echet
derivatives, and for every $z\in\mathbb{R}^n$
$$\mathcal{W}[z,z]=D^2v(x)[z,z]-D^2u(x)[z,z]\ \geq\ 0$$
by the hypothesis $D^2u(x)[w,w]\leq D^2v(x)[w,w]$ applied with $w=z$. Hence
\noderef{PosSemidefSymmQuadForm} applies to $\mathcal{W}$ and to the finite family
$s^i:=\sigma^i(x)$, $i=1,\dots,m$, and asserts that the matrix with entries
$\tfrac12\big(\mathcal{W}[s^j,s^i]+\mathcal{W}[s^i,s^j]\big)$ is positive semi-definite. By (2)
that matrix is $P$, so $P$ is positive semi-definite, which with Step 1 completes the proof.
\end{proof}


\subsection{\texorpdfstring{\texttt{TraceMonotone}}{TraceMonotone}}\label{node:TraceMonotone}
\begin{lemma}
Let $m\in\mathbb{N}$ and let $A,X,Y\in\mathbb{R}^{m\times m}$. Assume that $A$ is positive
semi-definite and that $X-Y$ is positive semi-definite, i.e.\ $Y\leq X$ in the order on symmetric
matrices of Section~2.1 of the paper. Then
$$\operatorname{Tr}(AY)\ \leq\ \operatorname{Tr}(AX).$$
\end{lemma}
\begin{proof}
Set $Z:=X-Y$, which is positive semi-definite by hypothesis. By linearity of the trace and
distributivity of the matrix product (preamble),
$$\operatorname{Tr}(AX)-\operatorname{Tr}(AY)=\operatorname{Tr}\big(A(X-Y)\big)
=\operatorname{Tr}(AZ),$$
so it suffices to prove $\operatorname{Tr}(AZ)\geq0$.

Since $A$ and $Z$ are positive semi-definite, each has a positive semi-definite square root
(preamble): there are symmetric $P,Q\in\mathbb{R}^{m\times m}$ with $P^2=A$ and $Q^2=Z$. Put
$B:=QP$, so that $B^T=P^TQ^T=PQ$. Using $\operatorname{Tr}(MN)=\operatorname{Tr}(NM)$ (preamble)
with $M=PPQ$ and $N=Q$ we get
$$\operatorname{Tr}(AZ)=\operatorname{Tr}(PPQQ)=\operatorname{Tr}(QPPQ)
=\operatorname{Tr}\big((QP)(PQ)\big)=\operatorname{Tr}\big(BB^T\big).$$
Finally, by the entrywise formula for the trace of a product (preamble),
$$\operatorname{Tr}\big(BB^T\big)=\sum_{i,k=1}^m B_{i,k}\,\big(B^T\big)_{k,i}
=\sum_{i,k=1}^m B_{i,k}^2\ \geq\ 0 .$$
Hence $\operatorname{Tr}(AX)-\operatorname{Tr}(AY)=\operatorname{Tr}(AZ)\geq0$, which is the
claim.
\end{proof}


\subsection{\texorpdfstring{\texttt{TraceMulSymmetrizeEq}}{TraceMulSymmetrizeEq}}\label{node:TraceMulSymmetrizeEq}
\begin{lemma}
Let $m\in\mathbb{N}$, let $A\in\mathbb{R}^{m\times m}$ be symmetric, i.e.\ $A^T=A$, and let
$B\in\mathbb{R}^{m\times m}$ be arbitrary. Then
$$\operatorname{Tr}\Big(A\cdot\tfrac12\big(B+B^T\big)\Big)=\operatorname{Tr}(AB).$$
Equivalently, replacing the second factor by its symmetrization does not change its trace against
a symmetric matrix. This is the matrix fact underlying the paper's remark, just after
\texttt{eq:maineq}, that because $a(\xi)$ is symmetric the equation
$\mathcal{L}u=-\operatorname{Tr}\big(a(\mathfrak{X}u)\mathfrak{X}\mathfrak{X}u\big)=0$ can also be
written with the symmetrized sub-elliptic second derivative
$\mathfrak{X}\mathfrak{X}^\star u$ in place of $\mathfrak{X}\mathfrak{X}u$.
\end{lemma}
\begin{proof}
First, $\operatorname{Tr}\big(AB^T\big)=\operatorname{Tr}(AB)$. Indeed, by the entrywise formula
for the trace of a product (preamble) and the symmetry of $A$,
$$\operatorname{Tr}\big(AB^T\big)=\sum_{i,k=1}^mA_{i,k}\big(B^T\big)_{k,i}
=\sum_{i,k=1}^mA_{i,k}B_{i,k}
=\sum_{i,k=1}^mA_{k,i}B_{i,k}=\operatorname{Tr}(AB),$$
where the first and last equalities are that formula applied to $AB^T$ and to $AB$, the second
uses $\big(B^T\big)_{k,i}=B_{i,k}$, and the third renames the summation indices
$i\leftrightarrow k$ and uses $A_{i,k}=A_{k,i}$, which is the hypothesis $A^T=A$.

Hence, using the linearity of the trace in the sense of the identities
$\operatorname{Tr}\big(A(B+C)\big)=\operatorname{Tr}(AB)+\operatorname{Tr}(AC)$ and
$\operatorname{Tr}\big(A(cB)\big)=c\operatorname{Tr}(AB)$ (preamble), with $C:=B^T$ and
$c:=\tfrac12$,
$$\operatorname{Tr}\Big(A\cdot\tfrac12\big(B+B^T\big)\Big)
=\tfrac12\operatorname{Tr}\big(A(B+B^T)\big)
=\tfrac12\operatorname{Tr}(AB)+\tfrac12\operatorname{Tr}\big(AB^T\big)
=\tfrac12\operatorname{Tr}(AB)+\tfrac12\operatorname{Tr}(AB)=\operatorname{Tr}(AB),$$
the third equality by the first paragraph.
\end{proof}


\subsection{\texorpdfstring{\texttt{LocalMaxOpIneq}}{LocalMaxOpIneq}}\label{node:LocalMaxOpIneq}
\begin{lemma}
Let $U\subset\mathbb{R}^n$ be open, let $\sigma=(\sigma^1,\dots,\sigma^m)$ be smooth vector fields
$\sigma^i:\mathbb{R}^n\to\mathbb{R}^n$ and let $a:\mathbb{R}^m\to\mathbb{R}^{m\times m}$ be such
that $a(\xi)$ is positive semi-definite (hence symmetric) for every
$\xi\in\mathbb{R}^m\setminus\{0\}$. Let $u,v:\mathbb{R}^n\to\mathbb{R}$ be of class $C^2$ on $U$
and let $x_0\in U$ be a local maximum point of $u-v$ with $\mathfrak{X}u(x_0)\neq0$. Then
$$\mathfrak{X}v(x_0)=\mathfrak{X}u(x_0),\qquad
\mathfrak{X}\mathfrak{X}^\star v(x_0)-\mathfrak{X}\mathfrak{X}^\star u(x_0)
\ \text{ is positive semi-definite},\qquad
\mathcal{L}v(x_0)\leq\mathcal{L}u(x_0),$$
the middle assertion being $\mathfrak{X}\mathfrak{X}^\star
u(x_0)\leq\mathfrak{X}\mathfrak{X}^\star v(x_0)$ in the order of Section~2.1 of the paper. This is
the first assertion of Lemma \texttt{lem:vweak} of the paper.
\end{lemma}
\begin{proof}
Since $U$ is open and $x_0\in U$, and $u,v$ are of class $C^2$ on $U$, both $u$ and $v$ are of
class $C^2$ in a neighbourhood of $x_0$ (namely on $U$ itself); in particular both are
differentiable at $x_0$. Each $\sigma^i$ is smooth, hence of class $C^1$ in a neighbourhood of
$x_0$.

\emph{Step 1: $\mathfrak{X}v(x_0)=\mathfrak{X}u(x_0)$.} Apply \noderef{HGradEqAtLocalMax} to
$\sigma$, $u$, $v$ and $x_0$: its hypotheses are that $u$ and $v$ are differentiable at $x_0$ and
that $x_0$ is a local maximum point of $u-v$, all of which hold. It gives
$\mathfrak{X}u(x_0)=\mathfrak{X}v(x_0)$. Write $\xi_0:=\mathfrak{X}u(x_0)=\mathfrak{X}v(x_0)$; by
hypothesis $\xi_0\neq0$.

\emph{Step 2: $Du(x_0)=Dv(x_0)$.} Apply \noderef{FDerivEqAtLocalMaxSub} to $u$, $v$ and $x_0$: its
hypotheses are that $u$ and $v$ are differentiable at $x_0$ and that $x_0$ is a local maximum point
of $u-v$, all of which hold. It gives $Du(x_0)=Dv(x_0)$.

\emph{Step 3: $D^2u(x_0)[w,w]\leq D^2v(x_0)[w,w]$ for every $w\in\mathbb{R}^n$.} Apply
\noderef{SndFDerivLeAtLocalMaxSub} to $U$, $u$, $v$ and $x_0$: its hypotheses are that $U$ is open,
that $x_0\in U$, that $u$ and $v$ are of class $C^2$ on $U$ and that $x_0$ is a local maximum point
of $u-v$, all of which hold by hypothesis. It gives exactly the claim.

\emph{Step 4: the symmetrized horizontal second derivatives.} Steps 2 and 3 are exactly the two
hypotheses of \noderef{HGradHHessComparison}, whose remaining hypotheses ($u,v$ of class $C^2$
near $x_0$, each $\sigma^i$ of class $C^1$ near $x_0$) also hold. It yields that
$$P:=\mathfrak{X}\mathfrak{X}^\star v(x_0)-\mathfrak{X}\mathfrak{X}^\star u(x_0)$$
is positive semi-definite, which is the second assertion.

\emph{Step 5: the trace comparison.} Put $A:=a(\xi_0)$. Since $\xi_0\neq0$, $A$ is positive
semi-definite by hypothesis; in particular $A$ is symmetric. Apply \noderef{TraceMonotone} with
this $A$, with $X:=\mathfrak{X}\mathfrak{X}^\star v(x_0)$ and
$Y:=\mathfrak{X}\mathfrak{X}^\star u(x_0)$: its hypotheses are that $A$ is positive semi-definite
(just checked) and that $X-Y=P$ is positive semi-definite (Step 4). It gives
\begin{equation*}\operatorname{Tr}\Big(A\,\mathfrak{X}\mathfrak{X}^\star u(x_0)\Big)\ \leq\
\operatorname{Tr}\Big(A\,\mathfrak{X}\mathfrak{X}^\star v(x_0)\Big). \tag{2}\end{equation*}

\emph{Step 6: symmetrization does not change the trace against a symmetric matrix.} The matrix
$A=a(\xi_0)$ is symmetric (Step 5), so \noderef{TraceMulSymmetrizeEq} applies to $A$ and to every
$B\in\mathbb{R}^{m\times m}$ and gives
$\operatorname{Tr}\big(A\cdot\tfrac12(B+B^T)\big)=\operatorname{Tr}(AB)$. Taking
$B=\mathfrak{X}\mathfrak{X}u(x_0)$ and then $B=\mathfrak{X}\mathfrak{X}v(x_0)$, and recalling that
by the definition of the symmetrization (\noderef{HHessSym}) one has
$\mathfrak{X}\mathfrak{X}^\star w(x_0)=\tfrac12\big(\mathfrak{X}\mathfrak{X}w(x_0)
+\mathfrak{X}\mathfrak{X}w(x_0)^T\big)$,
\begin{equation*}\operatorname{Tr}\Big(A\,\mathfrak{X}\mathfrak{X}^\star u(x_0)\Big)
=\operatorname{Tr}\Big(A\,\mathfrak{X}\mathfrak{X}u(x_0)\Big),\qquad
\operatorname{Tr}\Big(A\,\mathfrak{X}\mathfrak{X}^\star v(x_0)\Big)
=\operatorname{Tr}\Big(A\,\mathfrak{X}\mathfrak{X}v(x_0)\Big). \tag{3}\end{equation*}
This is the paper's remark, just after \texttt{eq:maineq}, that the equation can equivalently be
written with the symmetrized horizontal second derivative.

\emph{Step 7: conclusion.} By the definition of the operator (\noderef{SubellipticOp}) and
$\mathfrak{X}u(x_0)=\mathfrak{X}v(x_0)=\xi_0$,
$$\mathcal{L}u(x_0)=-\operatorname{Tr}\Big(A\,\mathfrak{X}\mathfrak{X}u(x_0)\Big),\qquad
\mathcal{L}v(x_0)=-\operatorname{Tr}\Big(A\,\mathfrak{X}\mathfrak{X}v(x_0)\Big).$$
Combining this with (3) and then (2),
$$-\mathcal{L}u(x_0)=\operatorname{Tr}\Big(A\,\mathfrak{X}\mathfrak{X}^\star u(x_0)\Big)
\ \leq\ \operatorname{Tr}\Big(A\,\mathfrak{X}\mathfrak{X}^\star v(x_0)\Big)=-\mathcal{L}v(x_0),$$
that is, $\mathcal{L}v(x_0)\leq\mathcal{L}u(x_0)$.
\end{proof}


\subsection{\texorpdfstring{\texttt{SndFDerivCompScalar}}{SndFDerivCompScalar}}\label{node:SndFDerivCompScalar}
\begin{lemma}
Let $n\in\mathbb{N}$, let $h\in C^2(\mathbb{R})$, let $w:\mathbb{R}^n\to\mathbb{R}$ and let
$x\in\mathbb{R}^n$. Assume that $w$ is of class $C^2$ in a neighbourhood of $x$. Then, for all
$z_1,z_2\in\mathbb{R}^n$,
$$D^2(h\circ w)(x)[z_1,z_2]=h''\big(w(x)\big)\,Dw(x)[z_1]\,Dw(x)[z_2]
+h'\big(w(x)\big)\,D^2w(x)[z_1,z_2].$$
\end{lemma}
\begin{proof}
Write $t:=h'(w(x))$ and $s:=h''(w(x))$, and let $V$ be an open neighbourhood of $x$ on which $w$ is
of class $C^2$.

\emph{Step 1: the first derivative of $h\circ w$ on $V$.} Since $h\in C^2(\mathbb{R})$, $h$ is
differentiable on $\mathbb{R}$; and $w$ is differentiable at every $y\in V$. Hence the chain rule
applies to $h\circ w$ at every $y\in V$ and gives
\begin{equation*}D(h\circ w)(y)=h'\big(w(y)\big)\,Dw(y)\qquad\text{for every }y\in V. \tag{1}\end{equation*}
In particular the map $y\mapsto D(h\circ w)(y)$ agrees on the neighbourhood $V$ of $x$ with the
product of the scalar function $y\mapsto h'(w(y))$ and the $\big(\mathbb{R}^n\big)^*$-valued
function $y\mapsto Dw(y)$.

\emph{Step 2: the two factors are differentiable at $x$.} Since $h\in C^2(\mathbb{R})$, the
function $h'$ is differentiable on $\mathbb{R}$, and $w$ is differentiable at $x$; so by the chain
rule $y\mapsto h'(w(y))$ is differentiable at $x$ with
$$D\big(y\mapsto h'(w(y))\big)(x)[z_1]=h''\big(w(x)\big)\,Dw(x)[z_1]=s\,Dw(x)[z_1]
\qquad\text{for every }z_1\in\mathbb{R}^n .$$
Since $w$ is of class $C^2$ on $V$, the map $y\mapsto Dw(y)$ is differentiable at $x$, with
derivative $z_1\mapsto D\big(y\mapsto Dw(y)\big)(x)[z_1]$.

\emph{Step 3: the product rule.} By Step 2 and the product rule for a scalar function times a
vector-valued map, the product of the two factors is differentiable at $x$, and by Step 1 it agrees
with $y\mapsto D(h\circ w)(y)$ on a neighbourhood of $x$; hence that map is differentiable at $x$
with the same derivative, which, using (1) at $y=x$ to evaluate the scalar factor, is
$$D\big(y\mapsto D(h\circ w)(y)\big)(x)[z_1]
=s\,\big(Dw(x)[z_1]\big)\,Dw(x)+t\,D\big(y\mapsto Dw(y)\big)(x)[z_1]$$
for every $z_1\in\mathbb{R}^n$. Evaluating this identity of linear functionals at
$z_2\in\mathbb{R}^n$ and using the definition of the second derivative from the preamble, i.e.\
$D^2g(x)[z_1,z_2]=\big(D(y\mapsto Dg(y))(x)[z_1]\big)[z_2]$,
$$D^2(h\circ w)(x)[z_1,z_2]=s\,Dw(x)[z_1]\,Dw(x)[z_2]+t\,D^2w(x)[z_1,z_2],$$
which is the claim.
\end{proof}


\subsection{\texorpdfstring{\texttt{HGradHHessComp}}{HGradHHessComp}}\label{node:HGradHHessComp}
\begin{lemma}
Let $\sigma=(\sigma^1,\dots,\sigma^m)$ be vector fields $\sigma^i:\mathbb{R}^n\to\mathbb{R}^n$, let
$h\in C^2(\mathbb{R})$, let $w:\mathbb{R}^n\to\mathbb{R}$ and let $x\in\mathbb{R}^n$. Assume that
$w$ is of class $C^2$ in a neighbourhood of $x$ and that every $\sigma^i$ is of class $C^1$ in a
neighbourhood of $x$. Then
$$\mathfrak{X}(h\circ w)(x)=h'\big(w(x)\big)\,\mathfrak{X}w(x)$$
and
$$\mathfrak{X}\mathfrak{X}(h\circ w)(x)=h'\big(w(x)\big)\,\mathfrak{X}\mathfrak{X}w(x)
+h''\big(w(x)\big)\,\mathfrak{X}w(x)\otimes\mathfrak{X}w(x),$$
where $\xi\otimes\eta\in\mathbb{R}^{m\times m}$ denotes the matrix with entries
$(\xi\otimes\eta)_{i,j}=\xi_i\,\eta_j$. These are the two identities
$\mathfrak{X}(h\circ w)=h'(w)\mathfrak{X}w$ and $\mathfrak{X}\mathfrak{X}(h\circ
w)=h'(w)\mathfrak{X}\mathfrak{X}w+h''(w)\,\mathfrak{X}w\otimes\mathfrak{X}w$ with which the proof
of Lemma \texttt{lem:tech} of the paper begins.
\end{lemma}
\begin{proof}
Write $t:=h'(w(x))$, $s:=h''(w(x))$ and $\xi:=\mathfrak{X}w(x)$.

\emph{Step 1: the first identity.} Let $V$ be an open neighbourhood of $x$ on which $w$ is of class
$C^2$. Since $h$ is differentiable on $\mathbb{R}$ and $w$ is differentiable on $V$, the chain
rule gives that $h\circ w$ is differentiable on $V$ with
\begin{equation*}D(h\circ w)(y)=h'\big(w(y)\big)\,Dw(y)\qquad\text{for every }y\in V. \tag{1}\end{equation*}
Evaluating (1) at $y=x$ and at the vector $\sigma^i(x)$, and using the definition of
\noderef{HGrad},
$$\big(\mathfrak{X}(h\circ w)(x)\big)_i=D(h\circ w)(x)\big[\sigma^i(x)\big]
=t\,Dw(x)\big[\sigma^i(x)\big]=t\,\xi_i,$$
for every $i$, which is the first identity.

\emph{Step 2: the second derivative of $h\circ w$.} Since $h\in C^2(\mathbb{R})$ and $w$ is of
class $C^2$ on $V$, the composition $h\circ w$ is of class $C^2$ on $V$, hence in a neighbourhood
of $x$. Moreover $h\in C^2(\mathbb{R})$ and $w$ is of class $C^2$ in a neighbourhood of $x$, which
are the hypotheses of \noderef{SndFDerivCompScalar}; applied to $h$, $w$ and $x$ it gives
\begin{equation*}D^2(h\circ w)(x)[z_1,z_2]=s\,Dw(x)[z_1]\,Dw(x)[z_2]+t\,D^2w(x)[z_1,z_2] \tag{2}\end{equation*}
for all $z_1,z_2\in\mathbb{R}^n$.

\emph{Step 3: the second identity.} Fix $i,j$. By Step 2 the function $h\circ w$ is of class $C^2$
in a neighbourhood of $x$, and $\sigma^i$ is of class $C^1$ in a neighbourhood of $x$; so
\noderef{HHessEntryFormula} applies to $h\circ w$ and gives
$$\big(\mathfrak{X}\mathfrak{X}(h\circ w)(x)\big)_{i,j}
=D^2(h\circ w)(x)\big[\sigma^j(x),\sigma^i(x)\big]
+D(h\circ w)(x)\big[D\sigma^i(x)\,\sigma^j(x)\big].$$
By (2) with $z_1=\sigma^j(x)$, $z_2=\sigma^i(x)$, and by the definition of \noderef{HGrad}, the
first summand equals
$$s\,Dw(x)\big[\sigma^j(x)\big]\,Dw(x)\big[\sigma^i(x)\big]
+t\,D^2w(x)\big[\sigma^j(x),\sigma^i(x)\big]
=s\,\xi_j\,\xi_i+t\,D^2w(x)\big[\sigma^j(x),\sigma^i(x)\big],$$
while by (1) the second summand equals
$t\,Dw(x)\big[D\sigma^i(x)\,\sigma^j(x)\big]$. Adding and applying
\noderef{HHessEntryFormula} to $w$ itself (which is of class $C^2$ near $x$),
$$\big(\mathfrak{X}\mathfrak{X}(h\circ w)(x)\big)_{i,j}
=t\Big(D^2w(x)\big[\sigma^j(x),\sigma^i(x)\big]+Dw(x)\big[D\sigma^i(x)\,\sigma^j(x)\big]\Big)
+s\,\xi_i\,\xi_j
=t\,\big(\mathfrak{X}\mathfrak{X}w(x)\big)_{i,j}+s\,\xi_i\,\xi_j .$$
Since $(\xi\otimes\xi)_{i,j}=\xi_i\xi_j$, this is precisely the second identity entrywise, and
hence as an identity of matrices.
\end{proof}


\subsection{\texorpdfstring{\texttt{PerturbationIdentity}}{PerturbationIdentity}}\label{node:PerturbationIdentity}
\begin{lemma}
Let $U\subset\mathbb{R}^n$ be open, let $\sigma=(\sigma^1,\dots,\sigma^m)$ be smooth vector fields
$\sigma^i:\mathbb{R}^n\to\mathbb{R}^n$, let $\beta\in\mathbb{R}$ and let
$a:\mathbb{R}^m\to\mathbb{R}^{m\times m}$ satisfy the homogeneity condition
$$a(t\xi)=t^\beta\,a(\xi)\qquad\text{for every }t>0\text{ and every }
\xi\in\mathbb{R}^m\setminus\{0\}.$$
Let $h\in C^2(\mathbb{R})$, let $w:\mathbb{R}^n\to\mathbb{R}$ be of class $C^2$ on $U$ and let
$x\in U$ be such that $\mathfrak{X}w(x)\neq0$ and $h'(w(x))>0$. Then
$$\mathcal{L}(h\circ w)(x)\ =\ -h'\big(w(x)\big)^{\beta}
\Big[h''\big(w(x)\big)\,\mathcal{E}\big(\mathfrak{X}w(x)\big)
-h'\big(w(x)\big)\,\mathcal{L}w(x)\Big],$$
where $\mathcal{E}$ is the ellipticity function of \noderef{EllipticityFn} and $t^\beta$ denotes
the real power of the positive real number $t=h'(w(x))$. This is Lemma \texttt{lem:tech}
(identity \texttt{eq:estLh}) of the paper.
\end{lemma}
\begin{proof}
Put $t:=h'(w(x))>0$, $s:=h''(w(x))$ and $\xi:=\mathfrak{X}w(x)\neq0$.

Since $U$ is open, $x\in U$ and $w$ is of class $C^2$ on $U$, the function $w$ is of class $C^2$ in
a neighbourhood of $x$; and each $\sigma^i$ is smooth, hence of class $C^1$ in a neighbourhood of
$x$.

\emph{Step 1: the two chain rules.} By \noderef{HGradHHessComp} (whose hypotheses are exactly
$h\in C^2(\mathbb{R})$, $w$ of class $C^2$ near $x$, and each $\sigma^i$ of class $C^1$ near $x$),
\begin{equation*}\mathfrak{X}(h\circ w)(x)=t\,\xi,\qquad
\mathfrak{X}\mathfrak{X}(h\circ w)(x)=t\,\mathfrak{X}\mathfrak{X}w(x)+s\,\xi\otimes\xi. \tag{1}\end{equation*}

\emph{Step 2: homogeneity of $a$ at the perturbed gradient.} Since $t>0$ and $\xi\neq0$, the
homogeneity hypothesis applies and gives
\begin{equation*}a\big(\mathfrak{X}(h\circ w)(x)\big)=a(t\,\xi)=t^{\beta}\,a(\xi). \tag{2}\end{equation*}

\emph{Step 3: expanding the operator.} By the definition of $\mathcal{L}$
(\noderef{SubellipticOp}) applied to $h\circ w$, then by (1), (2) and the linearity of the trace in
its second factor (preamble),
\begin{equation*}\mathcal{L}(h\circ w)(x)
=-\operatorname{Tr}\Big(a\big(\mathfrak{X}(h\circ w)(x)\big)\,
\mathfrak{X}\mathfrak{X}(h\circ w)(x)\Big)
=-t^{\beta}\Big[t\operatorname{Tr}\big(a(\xi)\,\mathfrak{X}\mathfrak{X}w(x)\big)
+s\operatorname{Tr}\big(a(\xi)\,\xi\otimes\xi\big)\Big]. \tag{3}\end{equation*}

\emph{Step 4: the two traces.} First, by the entrywise formula for the trace of a product
(preamble) and $(\xi\otimes\xi)_{k,i}=\xi_k\xi_i$,
$$\operatorname{Tr}\big(a(\xi)\,\xi\otimes\xi\big)=\sum_{i,k=1}^m
a_{i,k}(\xi)\,(\xi\otimes\xi)_{k,i}=\sum_{i,k=1}^m\xi_i\,a_{i,k}(\xi)\,\xi_k
=\big\langle a(\xi)\xi,\xi\big\rangle=\mathcal{E}(\xi),$$
by the definition of \noderef{EllipticityFn}. Second, since $\mathfrak{X}w(x)=\xi$, the definition
of $\mathcal{L}$ (\noderef{SubellipticOp}) applied to $w$ at $x$ reads
$\mathcal{L}w(x)=-\operatorname{Tr}\big(a(\xi)\,\mathfrak{X}\mathfrak{X}w(x)\big)$, i.e.
$$\operatorname{Tr}\big(a(\xi)\,\mathfrak{X}\mathfrak{X}w(x)\big)=-\mathcal{L}w(x).$$

\emph{Step 5: conclusion.} Substituting the two identities of Step 4 into (3),
$$\mathcal{L}(h\circ w)(x)=-t^{\beta}\big[-t\,\mathcal{L}w(x)+s\,\mathcal{E}(\xi)\big]
=-t^{\beta}\big[s\,\mathcal{E}(\xi)-t\,\mathcal{L}w(x)\big],$$
which, with $t=h'(w(x))$, $s=h''(w(x))$ and $\xi=\mathfrak{X}w(x)$, is the asserted identity.
\end{proof}


\subsection{\texorpdfstring{\texttt{MaxSetStability}}{MaxSetStability}}\label{node:MaxSetStability}
\begin{lemma}
Let $X$ be a topological space, let $K\subset X$ be non-empty and compact, let
$f:X\to\mathbb{R}$ be continuous on $K$ and let $W\subset X$ be open. Assume that every maximum
point of $f$ on $K$ lies in $W$, i.e.
$$x\in K\ \text{ and }\ f(y)\leq f(x)\ \ \forall y\in K\qquad\Longrightarrow\qquad x\in W .$$
Then there exists $\varepsilon>0$ with the following property: for every function
$g:X\to\mathbb{R}$ such that $|g(x)-f(x)|\leq\varepsilon$ for all $x\in K$, every maximum point of
$g$ on $K$ lies in $W$, i.e.
$$x\in K\ \text{ and }\ g(y)\leq g(x)\ \ \forall y\in K\qquad\Longrightarrow\qquad x\in W .$$
\end{lemma}
\begin{proof}
Since $K$ is non-empty and compact and $f$ is continuous on $K$, Weierstrass' theorem (preamble)
provides $x_0\in K$ with $f(y)\leq f(x_0)$ for every $y\in K$. Write $M_0:=f(x_0)=\max_Kf$.

\emph{Case 1: $K\setminus W=\emptyset$.} Then $K\subset W$, and $\varepsilon:=1$ works: for any
$g$ and any $x\in K$ we have $x\in W$ outright.

\emph{Case 2: $K\setminus W\neq\emptyset$.} Since $W$ is open, $K\setminus W=K\cap(X\setminus W)$
is the intersection of the compact set $K$ with a closed set, hence compact; and it is non-empty.
As $f$ is continuous on $K$ it is continuous on $K\setminus W\subset K$, so Weierstrass' theorem
(preamble) provides $x_1\in K\setminus W$ with
\begin{equation*}f(y)\leq f(x_1)\qquad\text{for every }y\in K\setminus W. \tag{1}\end{equation*}

We claim $f(x_1)<M_0$. Since $x_1\in K$ we have $f(x_1)\leq M_0$. If $f(x_1)=M_0$, then for every
$y\in K$ we get $f(y)\leq M_0=f(x_1)$, so $x_1$ is a maximum point of $f$ on $K$ and the hypothesis
forces $x_1\in W$, contradicting $x_1\in K\setminus W$. Hence $f(x_1)<M_0$.

Set
$$\varepsilon:=\frac{M_0-f(x_1)}{3}\ >\ 0 .$$
Let $g:X\to\mathbb{R}$ satisfy $|g(y)-f(y)|\leq\varepsilon$ for all $y\in K$, and let $x\in K$ be
such that $g(y)\leq g(x)$ for all $y\in K$. Suppose, for contradiction, that $x\notin W$. Then
$x\in K\setminus W$, so $f(x)\leq f(x_1)$ by (1), and therefore
$$g(x)\ \leq\ f(x)+\varepsilon\ \leq\ f(x_1)+\varepsilon .$$
On the other hand, taking $y=x_0\in K$ in the maximality of $x$ for $g$, and using
$|g(x_0)-f(x_0)|\leq\varepsilon$,
$$g(x)\ \geq\ g(x_0)\ \geq\ f(x_0)-\varepsilon\ =\ M_0-\varepsilon .$$
Combining the two displays, $M_0-\varepsilon\leq f(x_1)+\varepsilon$, i.e.
$$M_0-f(x_1)\ \leq\ 2\varepsilon\ =\ \tfrac{2}{3}\big(M_0-f(x_1)\big),$$
whence $M_0-f(x_1)\leq0$, contradicting $f(x_1)<M_0$. Therefore $x\in W$, which is the assertion.
\end{proof}


\subsection{\texorpdfstring{\texttt{HGradContinuousOn}}{HGradContinuousOn}}\label{node:HGradContinuousOn}
\begin{lemma}
Let $n,m\in\mathbb{N}$, let $\sigma=(\sigma^1,\dots,\sigma^m)$ be smooth vector fields
$\sigma^i:\mathbb{R}^n\to\mathbb{R}^n$, let $\Omega\subset\mathbb{R}^n$ be open and let
$u:\mathbb{R}^n\to\mathbb{R}$ be of class $C^2$ on $\Omega$. Then the horizontal gradient
$\mathfrak{X}u:\mathbb{R}^n\to\mathbb{R}^m$ is continuous on $\Omega$. This is the step ``by
assumption and continuity of $\mathfrak{X}u$'' at the start of the proof of Proposition
\texttt{lem:comp0} of the paper.
\end{lemma}
\begin{proof}
Since $u$ is of class $C^2$ on the open set $\Omega$, the differential map
$y\mapsto Du(y)$, with values in the space $\big(\mathbb{R}^n\big)^*$ of continuous linear
functionals on $\mathbb{R}^n$, is continuous on $\Omega$: a $C^2$ function on an open set is in
particular $C^1$ there, and for a function of class $C^1$ on an open set the differential is
continuous on that set (preamble).

Fix $i\in\{1,\dots,m\}$. By the definition of the horizontal gradient (\noderef{HGrad}), the
$i$-th component of $\mathfrak{X}u$ is
$$\big(\mathfrak{X}u(y)\big)_i=Du(y)\big[\sigma^i(y)\big].$$
The map $y\mapsto\sigma^i(y)$ is smooth on $\mathbb{R}^n$, hence continuous, and therefore
continuous on $\Omega$. The map $y\mapsto Du(y)$ is continuous on $\Omega$ by the previous
paragraph. The evaluation pairing
$$\big(\mathbb{R}^n\big)^*\times\mathbb{R}^n\to\mathbb{R},\qquad(\Lambda,\eta)\mapsto\Lambda[\eta],$$
is continuous, being bilinear and bounded on the finite-dimensional space
$\big(\mathbb{R}^n\big)^*\times\mathbb{R}^n$ (preamble). Composing the continuous map
$y\mapsto\big(Du(y),\sigma^i(y)\big)$ on $\Omega$ with this continuous pairing shows that
$y\mapsto Du(y)\big[\sigma^i(y)\big]$ is continuous on $\Omega$.

Thus every component of $\mathfrak{X}u$ is continuous on $\Omega$. A map into the
finite-dimensional space $\mathbb{R}^m$ is continuous on a set exactly when each of its $m$
components is (preamble), so $\mathfrak{X}u$ is continuous on $\Omega$.
\end{proof}


\subsection{\texorpdfstring{\texttt{QuadraticPerturbDeriv}}{QuadraticPerturbDeriv}}\label{node:QuadraticPerturbDeriv}
\begin{lemma}
Let $\lambda,u_0\in\mathbb{R}$ and let
$$h(s):=s+\lambda\,(s-u_0)^2\qquad(s\in\mathbb{R}).$$
Then $h\in C^2(\mathbb{R})$ and
$$h'(s)=1+2\lambda\,(s-u_0),\qquad h''(s)=2\lambda\qquad\text{for every }s\in\mathbb{R}.$$
These are the function $h_\lambda$ and the two derivative formulas displayed in the proof of
Proposition \texttt{lem:comp0} of the paper, where it is recorded that $h_\lambda'\geq1$ on the
range of $u$ (because $u-u_0\geq0$ there) and that $h_\lambda''=2\lambda$.
\end{lemma}
\begin{proof}
Write $q(s):=(s-u_0)^2$, so that $h(s)=s+\lambda\,q(s)$.

\emph{The derivative of $q$.} The affine map $s\mapsto s-u_0$ is differentiable at every
$s\in\mathbb{R}$ with derivative $1$, being the sum of the identity (derivative $1$) and the
constant $-u_0$ (derivative $0$) (preamble). Since $q(s)=(s-u_0)\cdot(s-u_0)$, the product rule
(preamble) gives, for every $s\in\mathbb{R}$,
$$q'(s)=1\cdot(s-u_0)+(s-u_0)\cdot1=2\,(s-u_0).$$

\emph{The derivative of $h$.} The identity has derivative $1$ at every point, and
$s\mapsto\lambda\,q(s)$ has derivative $\lambda\,q'(s)=\lambda\cdot2(s-u_0)$ by the previous
paragraph and the rule for a constant multiple (preamble). Hence, by the sum rule (preamble),
$h$ is differentiable at every $s\in\mathbb{R}$ with
$$h'(s)=1+\lambda\cdot2\,(s-u_0)=1+2\lambda\,(s-u_0),$$
which is the first asserted formula.

\emph{The second derivative of $h$.} By the previous paragraph the function $h'$ is given by
$h'(s)=1+2\lambda\,(s-u_0)$ for every $s\in\mathbb{R}$, i.e.\ $h'$ is the affine function
$s\mapsto1+2\lambda\,(s-u_0)$. Differentiating it at an arbitrary $s\in\mathbb{R}$: the map
$s\mapsto s-u_0$ has derivative $1$, so $s\mapsto2\lambda\,(s-u_0)$ has derivative
$2\lambda\cdot1=2\lambda$ by the constant-multiple rule, and adding the constant $1$ does not
change the derivative (preamble). Hence $h''(s)=2\lambda$ for every $s\in\mathbb{R}$, which is
the second asserted formula.

\emph{Regularity.} The function $h$ is a polynomial in $s$ of degree at most two, namely
$h(s)=\lambda s^2+(1-2\lambda u_0)s+\lambda u_0^2$. Polynomials are infinitely differentiable on
$\mathbb{R}$ (preamble), so in particular $h\in C^2(\mathbb{R})$.
\end{proof}


\section{Main results}\label{sec:main}
\subsection{\texorpdfstring{\texttt{ComparisonNondeg} (paper, Theorem 1.3)}{ComparisonNondeg}}\label{node:ComparisonNondeg}
\begin{theorem}
Let $n,m\in\mathbb{N}$, let $\sigma=(\sigma^1,\dots,\sigma^m)$ be smooth vector fields
$\sigma^i:\mathbb{R}^n\to\mathbb{R}^n$ satisfying the H\"ormander rank condition, i.e.\ such that
for every $x\in\mathbb{R}^n$ the linear span of $\{V(x):V\in\mathrm{HorFamily}(\sigma)\}$ is
$\mathbb{R}^n$ (\noderef{HormanderRank}), and let $X_i=\sigma^i(x)\cdot\nabla$. Let
$\beta\in\mathbb{R}$ and let $a:\mathbb{R}^m\to\mathbb{R}^{m\times m}$,
$a(\xi)=(a_{i,j}(\xi))_{i,j}$, be continuous on $\mathbb{R}^m\setminus\{0\}$ and such that, for
every $\xi\in\mathbb{R}^m\setminus\{0\}$, the matrix $a(\xi)$ is symmetric and positive
semi-definite,
$$\mathcal{E}(\xi)=\big\langle a(\xi)\xi,\xi\big\rangle>0,$$
and $a(t\xi)=t^\beta a(\xi)$ for every $t>0$. Let $\Omega\subset\mathbb{R}^n$ be a bounded open
set and let $u,v:\mathbb{R}^n\to\mathbb{R}$ be of class $C^2$ on $\Omega$ and continuous on
$\overline\Omega$. Assume that
$$\mathcal{L}u\leq0\leq\mathcal{L}v\quad\text{in }\Omega$$
in the classical sense, i.e.\ $u$ is a classical sub-solution and $v$ a classical super-solution of
$\mathcal{L}=0$ in $\Omega$ (\noderef{IsClassicalSubsol}, \noderef{IsClassicalSupersol}), and that
$u\leq v$ on $\partial\Omega$. If
$$\mathfrak{X}u(x)\neq0\qquad\text{for every }x\in\operatorname*{argmax}_\Omega(u-v)
:=\Big\{x\in\Omega:\ (u-v)(x)=\sup_\Omega(u-v)\Big\},$$
then $u\leq v$ in $\Omega$.
\end{theorem}
\begin{proof}
Suppose, for contradiction, that there is $\bar x\in\Omega$ with $u(\bar x)>v(\bar x)$.

Put $K:=\overline\Omega$ and $f:=u-v$. Since $\Omega$ is bounded, $K$ is compact (preamble), and
$K\neq\emptyset$ because $\bar x\in\Omega\subset K$. The function $f$ is continuous on $K$, since
$u$ and $v$ are. By Weierstrass' theorem (preamble) $f$ attains its maximum over $K$, and
\begin{equation*}M_0:=\max_Kf\ \geq\ f(\bar x)=u(\bar x)-v(\bar x)\ >\ 0. \tag{1}\end{equation*}

\emph{Step 1: every maximum point of $f$ on $K$ lies in $U:=\{x\in\Omega:\mathfrak{X}u(x)\neq0\}$.}
Let $x\in K$ be such that $f(y)\leq f(x)$ for all $y\in K$; then $f(x)=M_0>0$ by (1). Since
$K=\overline\Omega=\Omega\cup\partial\Omega$ (preamble), either $x\in\Omega$ or
$x\in\partial\Omega$; the latter is impossible because $f=u-v\leq0$ on $\partial\Omega$ by the
boundary hypothesis, while $f(x)>0$. Hence $x\in\Omega$. Moreover $\Omega\subset K$, so
$f(y)\leq f(x)$ for every $y\in\Omega$, i.e.\ $(u-v)(x)=\sup_\Omega(u-v)$ and
$x\in\operatorname*{argmax}_\Omega(u-v)$. The non-degeneracy hypothesis then gives
$\mathfrak{X}u(x)\neq0$, so $x\in U$.

\emph{Step 2: $U$ is open.} Since $u$ is of class $C^2$ on the open set $\Omega$, the map
$y\mapsto Du(y)$ is continuous on $\Omega$; each $\sigma^i$ is smooth, hence continuous. By the
definition of the horizontal gradient (\noderef{HGrad}), its $i$-th component is
$\big(\mathfrak{X}u(y)\big)_i=Du(y)\big[\sigma^i(y)\big]$; this is continuous on $\Omega$, being
the composition of the map $y\mapsto\big(Du(y),\sigma^i(y)\big)$, continuous on $\Omega$, with the
continuous evaluation pairing. As this holds for every $i\in\{1,\dots,m\}$, the map
$\mathfrak{X}u:\mathbb{R}^n\to\mathbb{R}^m$ is continuous on $\Omega$. (This is the paper's step
``by assumption and continuity of $\mathfrak{X}u$'' at the start of the proof of Proposition
\texttt{lem:comp0}.)

Now $\mathbb{R}^m\setminus\{0\}$ is open, and
$$U=\Omega\cap(\mathfrak{X}u)^{-1}\big(\mathbb{R}^m\setminus\{0\}\big)$$
is precisely the preimage of $\mathbb{R}^m\setminus\{0\}$ under the restriction of
$\mathfrak{X}u$ to $\Omega$. Since that restriction is continuous, $U$ is open in the subspace
topology of $\Omega$, i.e.\ $U=\Omega\cap V$ for some set $V$ open in $\mathbb{R}^n$. Since
$\Omega$ is itself open in $\mathbb{R}^n$, the intersection $U=\Omega\cap V$ is open in
$\mathbb{R}^n$.

\emph{Step 3: the perturbed sub-solution.} Let $u_0:=\min_Ku$ and $u_1:=\max_Ku$, which exist by
Weierstrass' theorem (preamble). For $\lambda>0$ set
$$h_\lambda(s):=s+\lambda(s-u_0)^2,\qquad u_\lambda:=h_\lambda\circ u .$$
Then $h_\lambda$ is a polynomial, hence $h_\lambda\in C^2(\mathbb{R})$, with
$$h_\lambda'(s)=1+2\lambda(s-u_0),\qquad h_\lambda''(s)=2\lambda\quad\text{for all }
s\in\mathbb{R}.$$
For $x\in K$ we have $u_0\leq u(x)\leq u_1$, hence $u(x)-u_0\geq0$ and therefore
\begin{equation*}h_\lambda'\big(u(x)\big)\geq1,\qquad
0\ \leq\ u_\lambda(x)-u(x)=\lambda\big(u(x)-u_0\big)^2\ \leq\ \lambda\,(u_1-u_0)^2
\qquad(x\in K). \tag{2}\end{equation*}
Moreover $u_\lambda=h_\lambda\circ u$ is of class $C^2$ on $\Omega$ (composition of the $C^\infty$
function $h_\lambda$ with the $C^2$ function $u$) and continuous on $K$.

\emph{Step 4: choice of $\lambda$ and of the perturbed maximum point.} The hypotheses of
\noderef{MaxSetStability} hold with the topological space $\mathbb{R}^n$, the non-empty compact
set $K$, the function $f=u-v$ (continuous on $K$) and the open set $W:=U$ (open by Step 2): by
Step 1, every maximum point of $f$ on $K$ lies in $U$. Hence there is $\varepsilon>0$ such that
\begin{equation*}\text{for every }g:\mathbb{R}^n\to\mathbb{R}\text{ with }|g-f|\leq\varepsilon\text{ on }K,
\text{ every maximum point of }g\text{ on }K\text{ lies in }U. \tag{3}\end{equation*}
Fix $\lambda>0$ with $\lambda(u_1-u_0)^2\leq\varepsilon$: if $u_1=u_0$ take $\lambda=1$, and
otherwise take $\lambda=\varepsilon/(u_1-u_0)^2>0$. Put $g:=u_\lambda-v$. By (2), for every
$x\in K$,
$$|g(x)-f(x)|=\big|u_\lambda(x)-u(x)\big|\leq\lambda(u_1-u_0)^2\leq\varepsilon .$$
Since $g$ is continuous on the non-empty compact set $K$, Weierstrass' theorem (preamble) gives
$x_\lambda\in K$ with $g(y)\leq g(x_\lambda)$ for every $y\in K$. By (3), $x_\lambda\in U$, that is
\begin{equation*}x_\lambda\in\Omega\qquad\text{and}\qquad\mathfrak{X}u(x_\lambda)\neq0. \tag{4}\end{equation*}

\emph{Step 5: $x_\lambda$ is a local maximum point of $u_\lambda-v$.} The set $\Omega$ is open and
contains $x_\lambda$ by (4), and $\Omega\subset K$, so $g(y)\leq g(x_\lambda)$ for every
$y\in\Omega$. Thus $\Omega$ is a neighbourhood of $x_\lambda$ on which
$u_\lambda-v\leq(u_\lambda-v)(x_\lambda)$, i.e.\ $x_\lambda$ is a local maximum point of
$u_\lambda-v$.

\emph{Step 6: $\mathcal{L}u_\lambda(x_\lambda)\geq0$.} Set $t:=h_\lambda'\big(u(x_\lambda)\big)$;
by (2) and $x_\lambda\in K$ we have $t\geq1>0$. Since $u$ is of class $C^2$ on the open set
$\Omega$ and $x_\lambda\in\Omega$, $u$ is of class $C^2$ in a neighbourhood of $x_\lambda$; each
$\sigma^i$ is smooth, hence of class $C^1$ near $x_\lambda$; and $h_\lambda\in C^2(\mathbb{R})$.
So the first identity of \noderef{HGradHHessComp} gives
\begin{equation*}\mathfrak{X}u_\lambda(x_\lambda)=t\,\mathfrak{X}u(x_\lambda)\ \neq\ 0, \tag{5}\end{equation*}
because $t\neq0$ and $\mathfrak{X}u(x_\lambda)\neq0$ by (4).

Now apply \noderef{LocalMaxOpIneq} with the open set $\Omega$, the smooth frame $\sigma$, the
matrix field $a$ (positive semi-definite away from $0$ by hypothesis), the functions $u_\lambda$
and $v$ (both of class $C^2$ on $\Omega$, by Step 3 and by hypothesis) and the point
$x_\lambda\in\Omega$, which is a local maximum point of $u_\lambda-v$ by Step 5 and satisfies
$\mathfrak{X}u_\lambda(x_\lambda)\neq0$ by (5). It gives
$$\mathfrak{X}v(x_\lambda)=\mathfrak{X}u_\lambda(x_\lambda)\neq0\qquad\text{and}\qquad
\mathcal{L}v(x_\lambda)\leq\mathcal{L}u_\lambda(x_\lambda).$$
Since $v$ is a classical super-solution in $\Omega$, $x_\lambda\in\Omega$ and
$\mathfrak{X}v(x_\lambda)\neq0$, the definition of \noderef{IsClassicalSupersol} gives
$\mathcal{L}v(x_\lambda)\geq0$. Therefore
\begin{equation*}\mathcal{L}u_\lambda(x_\lambda)\ \geq\ 0. \tag{6}\end{equation*}

\emph{Step 7: $\mathcal{L}u_\lambda(x_\lambda)<0$.} Put $\xi:=\mathfrak{X}u(x_\lambda)\neq0$.
Since $u$ is a classical sub-solution in $\Omega$, $x_\lambda\in\Omega$ and $\xi\neq0$, the
definition of \noderef{IsClassicalSubsol} gives
\begin{equation*}\mathcal{L}u(x_\lambda)\ \leq\ 0. \tag{7}\end{equation*}
Apply \noderef{PerturbationIdentity} with the open set $\Omega$, the smooth frame $\sigma$, the
matrix field $a$ and the exponent $\beta$ (the homogeneity hypothesis $a(t\xi)=t^\beta a(\xi)$ for
$t>0$, $\xi\neq0$ being part of our assumptions), with $h=h_\lambda\in C^2(\mathbb{R})$, $w=u$ of
class $C^2$ on $\Omega$, and $x=x_\lambda\in\Omega$, where
$\mathfrak{X}u(x_\lambda)=\xi\neq0$ and $h_\lambda'\big(u(x_\lambda)\big)=t\geq1>0$. Using
$h_\lambda''\equiv2\lambda$, it gives
\begin{equation*}\mathcal{L}u_\lambda(x_\lambda)=-t^{\beta}\Big[2\lambda\,\mathcal{E}(\xi)
-t\,\mathcal{L}u(x_\lambda)\Big]. \tag{8}\end{equation*}
We estimate the bracket from below. First, the function $\mathcal{E}$ appearing in the conclusion
of \noderef{PerturbationIdentity} is the ellipticity function of \noderef{EllipticityFn}, namely
$\mathcal{E}(\zeta)=\langle a(\zeta)\zeta,\zeta\rangle$, which is exactly the function displayed in
the ellipticity hypothesis of the present statement; since $\xi\neq0$, that hypothesis gives
$\mathcal{E}(\xi)>0$. Together with $\lambda>0$ this yields $2\lambda\,\mathcal{E}(\xi)>0$.
Second, $t>0$ and $\mathcal{L}u(x_\lambda)\leq0$ by (7), so
$-t\,\mathcal{L}u(x_\lambda)\geq0$. Hence
$$2\lambda\,\mathcal{E}(\xi)-t\,\mathcal{L}u(x_\lambda)\ \geq\ 2\lambda\,\mathcal{E}(\xi)\ >\ 0 .$$
Since $t>0$, the real power $t^{\beta}$ is positive. Therefore the right-hand side of (8) is the
product of the negative number $-t^\beta$ with a positive number, so
$$\mathcal{L}u_\lambda(x_\lambda)\ <\ 0,$$
which contradicts (6).

This contradiction shows that no $\bar x\in\Omega$ with $u(\bar x)>v(\bar x)$ exists, i.e.\
$u\leq v$ in $\Omega$.
\end{proof}


\subsection{\texorpdfstring{\texttt{CorNocritComparison} (paper, Corollary 1.4, comparison)}{CorNocritComparison}}\label{node:CorNocritComparison}
\begin{corollary}
Let $n,m\in\mathbb{N}$, let $\sigma=(\sigma^1,\dots,\sigma^m)$ be smooth vector fields
$\sigma^i:\mathbb{R}^n\to\mathbb{R}^n$ satisfying the H\"ormander rank condition, i.e.\ such that
for every $x\in\mathbb{R}^n$ the linear span of $\{V(x):V\in\mathrm{HorFamily}(\sigma)\}$ is
$\mathbb{R}^n$ (\noderef{HormanderRank}). Let $\beta\in\mathbb{R}$ and let
$a:\mathbb{R}^m\to\mathbb{R}^{m\times m}$ be continuous on $\mathbb{R}^m\setminus\{0\}$ and such
that, for every $\xi\in\mathbb{R}^m\setminus\{0\}$, the matrix $a(\xi)$ is symmetric and positive
semi-definite, $\mathcal{E}(\xi)=\langle a(\xi)\xi,\xi\rangle>0$, and $a(t\xi)=t^\beta a(\xi)$ for
every $t>0$. Let $\Omega\subset\mathbb{R}^n$ be a bounded open set and let
$u,v:\mathbb{R}^n\to\mathbb{R}$ be of class $C^2$ on $\Omega$ and continuous on $\overline\Omega$,
with $u$ a classical sub-solution and $v$ a classical super-solution of $\mathcal{L}=0$ in
$\Omega$ (\noderef{IsClassicalSubsol}, \noderef{IsClassicalSupersol}), and with $u\leq v$ on
$\partial\Omega$. If $v$, or $u$, has no horizontal critical point in $\Omega$, i.e.\ if
$$\mathfrak{X}v(x)\neq0\ \ \text{for every }x\in\Omega\qquad\text{or}\qquad
\mathfrak{X}u(x)\neq0\ \ \text{for every }x\in\Omega,$$
then $u\leq v$ in $\Omega$.
\end{corollary}
\begin{proof}
All hypotheses of \noderef{ComparisonNondeg} except its non-degeneracy hypothesis are among our
hypotheses verbatim (the data $\sigma,a,\beta,\Omega,u,v$, the smoothness of $\sigma$, the
H\"ormander rank condition, the continuity, symmetry, positive semi-definiteness, ellipticity and
homogeneity of $a$, the openness and boundedness of $\Omega$, the $C^2$ regularity of $u,v$ on
$\Omega$ and their continuity on $\overline\Omega$, the sub- and super-solution properties, and the
boundary inequality $u\leq v$ on $\partial\Omega$). So it suffices to verify its non-degeneracy
hypothesis:
\begin{equation*}\mathfrak{X}u(x)\neq0\qquad\text{for every }x\in\Omega\text{ with }(u-v)(x)=\sup_\Omega(u-v).
\tag{$\ast$}\end{equation*}

Let $x\in\Omega$ satisfy $(u-v)(y)\leq(u-v)(x)$ for every $y\in\Omega$. We distinguish the two
alternatives of the hypothesis.

\emph{Case 1: $\mathfrak{X}u(y)\neq0$ for every $y\in\Omega$.} Then in particular
$\mathfrak{X}u(x)\neq0$, since $x\in\Omega$.

\emph{Case 2: $\mathfrak{X}v(y)\neq0$ for every $y\in\Omega$.} Since $\Omega$ is open and
$x\in\Omega$, the set $\Omega$ is a neighbourhood of $x$ on which $u-v\leq(u-v)(x)$; hence $x$ is a
local maximum point of $u-v$. Moreover $u$ and $v$ are of class $C^2$ on the open set $\Omega$ and
$x\in\Omega$, so both are differentiable at $x$. Therefore \noderef{HGradEqAtLocalMax} applies and
gives
$$\mathfrak{X}u(x)=\mathfrak{X}v(x),$$
and the right-hand side is non-zero by the hypothesis of this case. Hence
$\mathfrak{X}u(x)\neq0$.

In both cases $\mathfrak{X}u(x)\neq0$, so $(\ast)$ holds. Applying
\noderef{ComparisonNondeg} we conclude that $u\leq v$ in $\Omega$.
\end{proof}


\subsection{\texorpdfstring{\texttt{CorNocritUniqueness} (paper, Corollary 1.4, uniqueness)}{CorNocritUniqueness}}\label{node:CorNocritUniqueness}
\begin{corollary}
Let $n,m\in\mathbb{N}$, let $\sigma=(\sigma^1,\dots,\sigma^m)$ be smooth vector fields
$\sigma^i:\mathbb{R}^n\to\mathbb{R}^n$ satisfying the H\"ormander rank condition, i.e.\ such that
for every $x\in\mathbb{R}^n$ the linear span of $\{V(x):V\in\mathrm{HorFamily}(\sigma)\}$ is
$\mathbb{R}^n$ (\noderef{HormanderRank}). Let $\beta\in\mathbb{R}$ and let
$a:\mathbb{R}^m\to\mathbb{R}^{m\times m}$ be continuous on $\mathbb{R}^m\setminus\{0\}$ and such
that, for every $\xi\in\mathbb{R}^m\setminus\{0\}$, the matrix $a(\xi)$ is symmetric and positive
semi-definite, $\mathcal{E}(\xi)=\langle a(\xi)\xi,\xi\rangle>0$, and $a(t\xi)=t^\beta a(\xi)$ for
every $t>0$. Let $\Omega\subset\mathbb{R}^n$ be a bounded open set and let
$u,v:\mathbb{R}^n\to\mathbb{R}$ be of class $C^2$ on $\Omega$ and continuous on $\overline\Omega$.
Assume that each of $u$ and $v$ is a solution of $\mathcal{L}=0$ in $\Omega$, i.e.\ that each of
them is simultaneously a classical sub-solution and a classical super-solution in $\Omega$
(\noderef{IsClassicalSubsol}, \noderef{IsClassicalSupersol}); that $u=v$ on $\partial\Omega$; and
that $v$ has no horizontal critical point in $\Omega$, i.e.\ $\mathfrak{X}v(x)\neq0$ for every
$x\in\Omega$. Then
$$u=v\qquad\text{on }\overline\Omega .$$
\end{corollary}
\begin{proof}
All the ambient hypotheses of \noderef{CorNocritComparison} on $\sigma$, $a$, $\beta$ and $\Omega$
(smoothness of $\sigma$, the H\"ormander rank condition, continuity, symmetry, positive
semi-definiteness, ellipticity and homogeneity of $a$, openness and boundedness of $\Omega$) are
among our hypotheses verbatim, and both $u$ and $v$ are of class $C^2$ on $\Omega$ and continuous
on $\overline\Omega$. We apply \noderef{CorNocritComparison} twice, exchanging the roles of the two
functions.

\emph{Step 1: $u\leq v$ in $\Omega$.} Apply \noderef{CorNocritComparison} with $u$ in the role of
the sub-solution and $v$ in the role of the super-solution. Indeed $u$ is a classical sub-solution
in $\Omega$ and $v$ is a classical super-solution in $\Omega$ (both being solutions); the boundary
hypothesis holds, since $u=v$ on $\partial\Omega$ gives $u\leq v$ on $\partial\Omega$; and the
disjunctive hypothesis holds through its first alternative, namely that the super-solution $v$ has
no horizontal critical point in $\Omega$, which is our hypothesis $\mathfrak{X}v(x)\neq0$ for every
$x\in\Omega$. The conclusion is $u(x)\leq v(x)$ for every $x\in\Omega$.

\emph{Step 2: $v\leq u$ in $\Omega$.} Apply \noderef{CorNocritComparison} again, now with $v$ in
the role of the sub-solution and $u$ in the role of the super-solution. Indeed $v$ is a classical
sub-solution in $\Omega$ and $u$ is a classical super-solution in $\Omega$ (both being solutions);
the boundary hypothesis holds, since $u=v$ on $\partial\Omega$ gives $v\leq u$ on $\partial\Omega$;
and the disjunctive hypothesis holds through its second alternative, which for this role assignment
reads that the sub-solution $v$ has no horizontal critical point in $\Omega$, again our hypothesis
$\mathfrak{X}v(x)\neq0$ for every $x\in\Omega$. The conclusion is $v(x)\leq u(x)$ for every
$x\in\Omega$.

\emph{Step 3: conclusion.} By Steps 1 and 2, $u(x)=v(x)$ for every $x\in\Omega$. By hypothesis
$u(x)=v(x)$ for every $x\in\partial\Omega$. Since
$\overline\Omega=\Omega\cup\partial\Omega$ (preamble), $u=v$ on $\overline\Omega$.
\end{proof}



\end{document}
